% Faithful Russian translation of the revised English manuscript, for author review.
\documentclass[runningheads]{llncs}
\usepackage[varvw]{newtxmath}
\usepackage[no-math]{fontspec}
\setmainfont{Liberation Serif}
\setsansfont{Liberation Sans}
\usepackage[english,russian]{babel}
\usepackage{graphicx,booktabs,amsmath,array,longtable}
\usepackage[hidelinks]{hyperref}
\hypersetup{pdftitle={Дополнительный экспериментальный протокол},pdfauthor={},pdfsubject={Русский перевод анонимной рукописи}}
\newcommand{\WIP}{\mathrm{WIP}}
\emergencystretch=2em
\begin{document}
\renewcommand{\keywordname}{{\bfseries Ключевые слова:}}
% Russian word hyphenation exceptions generated with Pyphen ru_RU.
\hyphenation{
аб-ля-ции
аб-ля-ций
аб-ля-ция
аб-со-лют-ная
аб-со-лют-ные
ав-то-ма-та
ав-то-ма-ти-че-ски
ав-то-ры
аг-ре-ги-ро-ван-ные
адап-та-цию
ак-ти-ва-ции
ак-ти-ва-ций
ак-ти-ва-цию
ак-ти-ва-ция
ак-ти-ви-ру-ет-ся
ак-тив-ные
ак-тив-ных
аль-тер-на-тив
аль-тер-на-ти-ва-ми
аль-тер-на-тив-ны-ми
аль-тер-на-тив-ных
аль-тер-на-ти-вы
ана-лиз
ана-ли-за
ана-ли-зи-ру-ют-ся
ана-ли-зы
ана-ли-ти-че-ский
ано-ним-ная
ано-ним-ной
ано-ним-ные
ар-те-фак-там
ар-хив
ар-хи-ве
ар-хи-ви-ро-ван-ной
ар-хи-вы
атри-бу-та
атри-бу-там
атри-бу-тах
атри-бу-ты
аудит
ауди-те
ба-за-ми
ба-зе
ба-зо-вая
ба-зо-во-го
ба-зо-вой
ба-зо-вые
ба-зо-вый
ба-зы
без-опас-ные
без-услов-но
без-услов-но-го
без-услов-ным
биб-лио-те-ка
биб-лио-те-ке
биб-лио-те-ку
биз-нес
би-нар-ная
би-нар-но-го
би-нар-ные
би-нар-ных
бла-го-при-ят-но
бла-го-при-ят-ной
бли-жай-шие
блок
бло-ка
бло-ки
бло-ков
блоч-ные
бо-га-той
бо-лее
боль-шей
бон-фер-ро-ни
бу-ду-ще-го
бу-ду-щее
бу-ду-щем
бу-ду-ще-му
бу-ду-щие
бу-ду-щих
бу-стинг
бу-стин-га
бут-стреп
бы-ли
быть
важ-ную
ва-ли-да-ции
ва-ли-да-ци-он-ной
ва-ли-да-ци-он-ные
ва-ли-да-ци-он-ный
ва-ли-да-ци-он-ным
ва-ли-да-ци-он-ных
ва-ли-да-ция
ва-ри-ант
ва-ри-ан-та
ва-ри-ан-та-ми
ва-ри-ан-тах
ва-ри-ан-те
ва-ри-ан-ты
ва-рьи-ру-ет
ва-рьи-руя
вве-де-ние
вверх
вдвое
ве-ли-чи-на
ве-ли-чи-ну
вер-ну-лась
ве-ро-ят-но-стей
ве-ро-ят-но-сти
ве-ро-ят-ност-ной
ве-ро-ят-ность
ве-ро-ят-но-стя-ми
вер-сии
вер-си-о-ни-ро-ван-ном
вер-сию
ве-са
ве-са-ми
ве-сах
вет-ви
вза-и-мо-дей-ствие
вза-и-мо-дей-ствия
вза-и-мо-за-ме-ня-е-мых
ви-ди-мые
ви-ди-мый
ви-ди-мых
вклад
вклю-ча-ет
вклю-ча-ет-ся
вклю-чать
вклю-ча-ют
вклю-ча-ют-ся
вклю-чая
вклю-че-ние
вклю-че-ний
вклю-че-ния
вклю-че-ны
вла-де-лец
вли-я-ет
вли-я-ние
вли-я-нии
вли-я-ния
вли-я-ют
вло-жен-ные
вло-жен-ным
вме-ша-тельств
вме-ша-тель-ства
вме-ша-тель-ствах
вме-ша-тель-ство
вме-ша-тель-ством
вне-дре-ние
вне-дре-нии
вне-дре-ния
внеш-них
вно-сить
вну-трен-ней
вну-три
вну-три-про-цесс-ные
воз-врат
воз-вра-та
воз-вра-там
воз-вра-та-ми
воз-вра-тах
воз-вра-те
воз-вра-тов
воз-вра-том
воз-вра-ты
воз-вра-ща-е-мую
воз-вра-ще-ния
воз-мож-ное
воз-мож-ную
воз-мож-ны
воз-мож-ные
воз-мож-ным
воз-му-ще-ние
воз-му-ще-нию
воз-об-нов-ле-ния
во-круг
во-прос
во-про-сы
во-семь
во-се-мью
вос-про-из-ве-де-ния
вос-про-из-во-ди-мость
вос-про-из-во-дит
вос-про-из-во-дят
вос-ста-нав-ли-ва-ет
вос-ста-нав-ли-ва-ет-ся
вос-ста-нав-ли-ва-ют
вос-ста-нов-ле-ние
вос-ста-нов-ле-ни-ем
вос-ста-нов-лен-ное
вос-ста-нов-лен-ном
вось-ми
вред-ное
вре-ме-на
вре-ме-ни
вре-мен-ное
вре-мен-ный
вре-мя
все-гда
все-го
всей
всем
всех
всплеск
всплес-ка
всплес-ках
всплес-ке
всплес-ки
всплес-ков
всплес-ком
вспо-мо-га-тель-ные
вспо-мо-га-тель-ных
вто-рич-ное
вто-рич-ны
вто-рич-ных
вто-рой
вто-ром
вхо-да-ми
вхо-дит
вход-ные
вхо-дов
вхо-ды
вы-би-ра-ет
вы-би-ра-ет-ся
вы-би-ра-ют
вы-би-ра-ют-ся
вы-би-рая
вы-бор
вы-бо-ра
вы-бо-ра-ми
вы-бо-ре
вы-бор-ка
вы-бор-кам
вы-бор-ках
вы-бор-ке
вы-бор-ки
вы-бор-кой
вы-бо-рок
вы-бо-ром
вы-бо-ру
вы-бран
вы-бран-ная
вы-бран-но-го
вы-бран-ное
вы-бран-ной
вы-бран-ную
вы-бран-ные
вы-бран-ным
вы-бран-ных
вы-бра-но
вы-брать
вы-во-да
вы-во-дит
вы-во-дов
вы-во-ды
вы-во-дят
вы-де-ле-ни-ем
вы-де-лить
вы-де-ля-ет
вы-зван-ный
вы-игрыш
вы-игры-ша
вы-пол-нен
вы-пол-не-ние
вы-пол-нен-ных
вы-пол-ня-ет
вы-пол-ня-ют-ся
вы-рав-ни-ва-ни-ям
вы-со-кая
вы-со-кой
вы-ход
вы-хо-де
вы-хо-дит
вы-ход-ные
вы-хо-дом
вы-чис-ле-ний
вы-чис-ли-тель-но-го
вы-чис-ли-тель-ную
вы-чис-ля-ет-ся
вы-ше
вы-явил
вы-яви-ла
вы-яв-ле-ние
вы-яв-ле-ния
вы-яв-лен-ная
вы-яв-лен-ной
вы-яв-лен-ную
вы-яв-лен-ные
вы-яв-лен-ны-ми
вы-яв-лен-ных
вы-яв-ля-ет
вы-яв-лять
вы-яв-ля-ют
га-ран-ти-ру-ет
га-ран-тия
ге-не-ра-тив-но-го
ге-не-ра-тив-ный
ге-не-ра-тив-ных
ге-не-ра-тор
ге-не-ра-то-ра
ге-не-ра-то-ром
ге-не-ра-то-ру
ге-не-ра-то-ры
ге-не-ра-ции
ги-брид-ной
ги-по-те-ти-че-ское
ги-сто-грамм-ный
глав-ных
гло-баль-ная
гло-баль-но
гло-баль-ной
гло-баль-ный
глу-бо-ко-го
го-ри-зонт
го-ри-зон-та
го-ри-зон-таль
го-ри-зон-таль-ная
го-ри-зон-тах
го-ри-зон-те
го-ри-зон-том
го-ри-зон-ту
го-то-вит
го-тов-но-сти
го-то-вую
го-то-вые
го-то-вых
гра-ди-ент-ный
гра-ни-ца
гра-ни-ца-ми
гра-ни-цей
гра-ни-цы
гра-фи-ков
гру-бую
груп-пам
груп-пе
груп-пы
да-вал
да-ва-ло
даль-ней-шую
дан-ной
дан-ном
дан-ные
дан-ным
дан-ных
да-ют
даёт
два-жды
две-на-дцать
дву-мя
дву-сто-рон-ние
дву-сто-рон-ни-ми
двух
дей-ствие
дей-ствий
дей-ствия
дей-стви-ям
дей-стви-я-ми
дей-стви-ях
дей-ству-ет
де-ла-ют
де-лит-ся
де-лят
де-лят-ся
делён-ная
делён-но-му
делён-ные
де-неж-ные
день-гах
де-ре-ва
де-ре-во
де-тер-ми-ни-ро-ван-ные
де-тер-ми-ни-ро-ван-ных
диа-гно-сти-ка
диа-гно-сти-ке
диа-гно-сти-кой
диа-гно-сти-ку
диа-гно-сти-ру-ет
диа-гно-сти-че-ски
диа-гно-сти-че-ские
диа-гно-сти-че-ский
диа-гно-сти-че-ских
диа-гно-сти-че-ско-го
диа-па-зон
диа-па-зо-не
диа-па-зо-ном
диа-па-зо-ну
ди-зайн
дис-пер-сии
дис-пер-сию
дли-на
дли-ны
дли-тель-но-го
дли-тель-ной
дли-тель-но-стей
дли-тель-но-сти
дли-тель-ность
дли-тель-но-стях
дли-тель-ную
дли-тель-ных
до-бав-ле-ние
до-бав-ле-ни-ем
до-бав-ле-нии
до-бав-ле-ния
до-бав-лен-ные
до-бав-ле-ны
до-бав-ля-ет
до-бав-ля-ют-ся
до-ве-ри-тель-ный
до-ве-ри-тель-ных
до-ве-рия
до-ка-зы-ва-ет
до-ка-зы-ва-ют
до-кри-зис-ной
до-ку-мен-ти-ро-ван-ны-ми
дол-жен
долж-на
долж-но
до-ли
доль-ше
до-лю
до-ми-ни-ру-ю-щая
до-пол-ни-тель-ная
до-пол-ни-тель-но-го
до-пол-ни-тель-ное
до-пол-ни-тель-ную
до-пол-ни-тель-ные
до-пол-ни-тель-ный
до-пол-ни-тель-ных
до-пол-ня-ют
до-пуск
до-пус-ка
до-пус-ка-ет
до-пус-ка-ют
до-пус-ка-ют-ся
до-пус-ка-ю-щих
до-пу-сти-мо
до-пу-сти-мое
до-пу-сти-мы
до-пу-сти-мые
до-пу-сти-мых
до-пу-ще-ни-ем
до-ста-точ-ной
до-сти-га-ет
до-стиг-ших
до-сти-жи-мым
до-сти-жи-мых
до-ступ-но
до-ступ-ное
до-ступ-но-сти
до-ступ-ность
до-ступ-ны
до-ступ-ные
до-сту-пом
дру-гие
дру-гих
дру-го-го
дру-гой
еди-ниц
еди-ни-ца
еди-ни-ца-ми
еди-ни-цей
еди-ни-цу
еди-ни-цы
един-ствен-ный
если
есть
жиз-нен-но-го
жиз-нен-ном
жур-нал
жур-на-ла
жур-на-лам
жур-на-лах
жур-на-лов
жур-на-лы
за-вер-ша-ют-ся
за-вер-ше-ние
за-вер-ше-ни-ем
за-вер-ше-нии
за-вер-ше-ния
за-вер-ше-ны
за-вер-шён-ных
за-ви-си-мо-стей
за-ви-си-мо-сти
за-ви-си-мость
за-ви-си-мо-стя-ми
за-ви-си-мо-стях
за-ви-си-мые
за-ви-си-мым
за-ви-си-мы-ми
за-ви-сит
за-вы-шать
за-да-вать
за-дан
за-дан-ная
за-дан-но-го
за-дан-ное
за-дан-ной
за-дан-ную
за-дан-ные
за-дан-ным
за-дан-ны-ми
за-дан-ных
за-да-ны
за-дать
за-дач
за-да-ча
за-да-чам
за-да-ча-ми
за-да-чах
за-да-чей
за-да-чи
за-да-чу
за-да-ют
за-да-ют-ся
за-даёт
за-даёт-ся
за-держ-ку
за-кан-чи-ва-ет
за-кан-чи-ва-ют-ся
за-клю-че-ние
за-клю-чи-тель-ный
за-ко-на
за-мед-ле-ние
за-мед-лять
за-мед-ля-ют-ся
за-ме-ня-ет-ся
за-ме-ня-ют-ся
за-ни-ма-ют
за-но-во
за-ня-тость
за-ня-то-стью
за-ня-тых
за-пас
за-па-са
за-пи-сан-ным
за-пи-са-ны
за-пи-си
за-пись
за-пуск
за-пус-ка-ет
за-пус-ке
за-пус-ков
за-ра-нее
за-тем
за-то-ра
за-тяж-но-го
за-тяж-ной
за-фик-си-ро-ва-ли
за-фик-си-ро-ван-ную
за-фик-си-ро-ван-ные
за-фик-си-ро-ван-ный
за-фик-си-ро-ван-ны-ми
за-фик-си-ро-ван-ных
за-фик-си-ро-ва-ны
за-яв-ка
за-яв-ля-ем
за-яв-ля-ют
здесь
зер-на
зер-но
зна-ко-мых
зна-ме-на-тель
зна-ме-на-те-ля
зна-ния
зна-че-ние
зна-че-нию
зна-че-ния
иг-но-ри-ро-ва-ла
иден-ти-фи-ка-тор
иден-ти-фи-ка-то-рам
иден-ти-фи-ка-то-ры
иден-ти-фи-ка-ции
иден-ти-фи-ци-ру-ют
иден-тич-ны
идут
из-вест-ная
из-вест-ной
из-вест-но-му
из-вест-ную
из-вест-ный
из-вест-ным
из-ме-не-ние
из-ме-не-нии
из-ме-не-ния
из-ме-ня-ет
из-ме-ня-ют
из-ме-няя
из-менён-ным
из-ме-ре-ние
из-ме-ре-ний
из-ме-ре-ния
из-ме-ря-ет
из-ме-ря-ет-ся
из-ме-ря-ют
из-ме-ря-ют-ся
изу-ча-ем
изу-ча-ет-ся
име-ет
имен-но
име-ют
ими-та-ции
ими-та-ци-он-ная
ими-та-ци-он-но-го
ими-та-ци-он-ное
ими-та-ци-он-ную
ими-та-ци-он-ные
ими-та-ция
им-пуль-са
им-пуль-сах
им-пульс-ные
им-пульс-ных
ина-че
ин-дек-си-ро-ван-ные
ин-же-нер
ин-же-не-рию
ин-же-нер-ная
ин-же-нер-ное
ин-же-нер-ные
ин-же-нер-ны-ми
ин-же-не-ром
ин-же-не-ры
ин-стру-мент
ин-стру-мен-та
ин-те-грал
ин-те-гра-лы
ин-те-граль-ной
ин-тен-сив-но-сти
ин-тен-сив-ность
ин-тен-сив-но-стью
ин-тен-сив-но-стя-ми
ин-тер-вал
ин-тер-ва-ла
ин-тер-ва-ла-ми
ин-тер-ва-лом
ин-тер-ва-лы
ин-те-ре-сов
ин-тер-пре-та-цию
ин-тер-пре-ти-ру-е-мые
ин-тер-пре-ти-ру-ет-ся
ин-фор-ма-ции
ин-фор-ма-ци-он-ной
ин-фор-ма-ци-он-ных
ин-фор-ма-цию
ин-ци-ден-тов
иных
ис-клю-ча-ет
ис-клю-чать
ис-клю-ча-ют
ис-клю-че-ны
ис-клю-чён
ис-клю-чён-но-го
ис-клю-чён-ные
ис-пол-не-ние
ис-пол-не-ния
ис-пол-нен-ная
ис-пол-ни-мая
ис-пол-ни-мой
ис-пол-ни-мую
ис-пол-ни-мых
ис-пол-ни-те-лей
ис-пол-ни-те-лем
ис-пол-ни-тель
ис-пол-ни-те-ля
ис-пол-ня-е-мые
ис-пол-ня-ет
ис-поль-зо-вал
ис-поль-зо-вал-ся
ис-поль-зо-ва-ние
ис-поль-зо-ва-ни-ем
ис-поль-зо-ван-ные
ис-поль-зо-ван-ных
ис-поль-зу-ем
ис-поль-зу-ет
ис-поль-зу-ет-ся
ис-поль-зу-ют
ис-поль-зу-ют-ся
ис-поль-зуя
ис-прав-ле-ни-ем
ис-прав-ле-ния
ис-прав-лен-ный
ис-прав-лен-ным
ис-пы-та-ние
ис-сле-до-ва-ние
ис-сле-до-ва-ний
ис-сле-до-ва-ния
ис-сле-до-ва-тель-ские
ис-сле-до-ва-тель-ский
ис-сле-до-ва-тель-ском
ис-сле-ду-ем
ис-тин-ная
ис-тин-но-го
ис-тин-ное
ис-тин-ной
ис-тин-ном
ис-тин-но-му
ис-тин-ные
ис-тин-ный
ис-тин-ным
ис-ти-ны
ис-то-рии
ис-то-ри-че-ские
ис-то-ри-че-ским
ис-то-рия
ис-то-ри-ях
ис-точ-ник
ис-точ-ни-ки
ис-ход
ис-хо-дам
ис-хо-да-ми
ис-ход-ная
ис-ход-но-го
ис-ход-ное
ис-ход-ной
ис-ход-ном
ис-ход-но-му
ис-ход-ные
ис-ход-ный
ис-ход-ным
ис-ход-ны-ми
ис-ход-ных
ис-хо-ды
ис-че-за-ет
ис-чер-пы-ва-ю-ще
каж-дая
каж-до-го
каж-дой
каж-дом
каж-до-му
каж-дую
каж-дые
каж-дый
ка-кие
ка-лен-дар-ные
ка-лен-дарь
ка-либ-ро-ван-но-го
ка-либ-ров-ка
ка-либ-ров-ке
ка-либ-ров-ки
ка-либ-ро-воч-ные
ка-либ-ро-воч-ный
ка-либ-ро-воч-ных
ка-либ-ру-ет
кам-па-нии
кан-ди-дат
кан-ди-дат-ные
кан-ди-дат-ных
кан-ди-да-тов
кан-ди-да-том
кан-ди-да-ты
кар-точ-ка
кар-точ-ки
кар-точ-ку
ка-та-лог
ка-та-ло-га
ка-та-ло-ге
ка-че-ства
ка-че-стве
ка-че-ством
квад-ра-тич-ной
кви-тан-ции
класс
клас-са
клас-си-фи-ка-тор
клас-си-фи-ка-то-ра
клас-сов
клас-су
клю-че-вые
ко-гор-ты
ко-да
ком-би-на-ций
ком-пен-си-ру-ет
ком-пи-ля-тор
ком-пи-ля-то-ра
ком-по-зи-ци-он-ный
ком-по-нент
ком-по-нен-та
ком-по-нен-тов
ком-про-мисс
ком-про-мис-сы
кон-вей-ер
ко-неч-но-го
ко-неч-ной
ко-неч-ном
ко-неч-ную
ко-неч-ные
ко-неч-ный
кон-ку-рен-ции
кон-ку-ри-ру-ю-щи-ми
кон-ку-ри-ру-ю-щих
кон-текст
кон-текст-ный
кон-тро-ле
кон-тро-ли-ру-е-мой
кон-тро-ли-ру-е-мую
кон-тро-ли-ру-е-мые
кон-тро-ли-ру-е-мый
кон-трол-лер
кон-троль-ном
кон-троль-ные
кон-троль-ный
кон-троль-ным
контр-фак-ти-че-ский
контр-фак-ти-че-ско-го
кон-фи-гу-ра-ции
кон-флик-тов
кон-ца
ко-ор-ди-нат
ко-ор-ди-на-ты
кор-рек-ти-ру-ют
кор-рект-ная
ко-то-рая
ко-то-рое
ко-то-ром
ко-то-рые
ко-то-рый
ко-то-рым
ко-эф-фи-ци-ент
ко-эф-фи-ци-ен-тов
ко-эф-фи-ци-ен-ту
ко-эф-фи-ци-ен-ты
крат-кие
крат-кое
крат-ко-сроч-ный
кри-зис
кри-зи-са
кри-зис-ные
кри-зи-сов
кри-те-рий
кри-те-рию
кри-те-рия
кри-те-ри-я-ми
кросс
ле-жат
лек-си-че-ским
ли-бо
ли-ния
ли-те-ра-ту-ры
лишь
ло-га-рифм
ло-га-риф-ма
ло-га-риф-мам
ло-га-риф-ми-че-ское
ло-га-риф-мы
ло-ги-сти-че-ская
ло-гли-ней-ная
ло-гнор-маль-но-го
ло-гнор-маль-ный
лож-ной
лож-ные
ло-ка-ли-зу-ет
ло-каль-ной
ло-каль-но-му
ло-каль-ный
луч-шее
лю-бо-го
лю-бой
мак-си-маль-но-го
мак-си-маль-ный
мак-си-ми-зи-ру-ет-ся
мак-си-мум
мак-си-му-мы
мар-ки-ров-ки
мар-ков-ских
марш-ру-та-ми
марш-ру-тах
марш-ру-ти-за-ции
марш-ру-ти-за-ция
марш-ру-тов
марш-ру-ты
мас-си-вов
мас-штаб
мас-шта-би-ру-е-мо-сти
мас-шта-би-ру-ют-ся
мас-шта-бы
ма-те-ри-а-лах
ма-те-ри-а-лов
ма-те-ри-а-лы
ма-шин-но-го
ма-ши-но-чи-та-е-мую
меж-ду
ме-нее
мень-ше
мень-шую
ме-ня-ет
ме-ня-ет-ся
ме-нять
ме-нять-ся
ме-ня-ют
ме-ня-ют-ся
ме-няя
мест
ме-ста
ме-та-ста-биль-ных
мет-ка
мет-ка-ми
мет-ки
ме-тод
ме-то-да
ме-то-ди-ка
ме-то-ди-ке
ме-то-ди-ки
ме-то-дов
ме-то-дом
ме-то-ды
ме-ток
ме-ха-низм
ме-ха-низ-ма
ме-ха-низ-мам
ме-ха-низ-мах
ме-ха-низ-ме
ме-ха-низ-мов
ме-ха-низ-мом
ме-ха-низ-му
ме-ха-низ-мы
ми-ни-маль-ном
ми-ни-маль-ную
ми-ни-ми-за-ция
ми-ни-мум
ми-нус
ми-ра
ми-рах
ми-ре
ми-ров
мно-го-кри-те-ри-аль-ное
мно-го-кри-те-ри-аль-ные
мно-же-ствен-ные
мно-жи-те-ли
мно-жи-тель
мно-жи-те-ля
мо-гут
мо-де-лей
мо-де-ли
мо-де-ли-ро-ва-ние
мо-дель
мо-де-лью
мо-де-ля-ми
мо-жет
мож-но
мо-мент
мо-мен-та
мо-мен-ты
мо-ни-то-ринг
мо-ни-то-рин-га
мо-ти-ви-ро-ва-ло
мощ-но-сти
мощ-ность
муль-ти-но-ми-аль-ная
на-блю-дав-ше-го-ся
на-блю-дав-шей-ся
на-блю-дав-шем-ся
на-блю-дав-шим-ся
на-блю-дав-ших-ся
на-блю-да-е-мая
на-блю-да-е-мо-го
на-блю-да-е-мое
на-блю-да-е-мом
на-блю-да-е-мо-му
на-блю-да-е-мую
на-блю-да-е-мые
на-блю-да-е-мый
на-блю-да-е-мым
на-блю-да-е-мы-ми
на-блю-да-е-мых
на-блю-да-ет
на-блю-да-тель-ная
на-блю-да-тель-ную
на-блю-да-тель-ные
на-блю-да-ют-ся
на-блю-де-ние
на-блю-де-ний
на-блю-де-ния
на-блю-де-ни-ям
на-бор
на-бо-ра
на-бо-рам
на-бо-ра-ми
на-бо-рах
на-бо-ре
на-бо-ров
на-бо-ром
на-бо-ры
на-груз-ка
на-груз-ка-ми
на-груз-ках
на-груз-ке
на-груз-ки
на-груз-кой
на-груз-ку
на-гру-зок
на-гру-зоч-ная
на-гру-зоч-но-го
на-гру-зоч-ное
на-гру-зоч-ную
на-гру-зоч-ные
на-гру-зоч-ный
на-гру-зоч-ным
на-гру-зоч-ны-ми
на-гру-зоч-ных
надёж-ность
на-зна-че-ние
на-зна-че-ния
наи-боль-шая
наи-боль-шей
наи-мень-шее
на-ко-нец
на-коп-ле-ние
на-коп-лен-ной
на-ли-чие
на-о-бо-рот
на-пи-са-нии
на-при-мер
на-про-тив
на-ру-ша-ет
на-ру-ша-ю-щие
на-ру-шен
на-ру-ше-ний
на-ру-ше-ния
на-ря-ду
на-стра-и-ва-е-мые
на-строй-ке
на-строй-ки
на-сы-ща-ю-щей-ся
на-хо-дят-ся
на-хо-жде-ния
на-ча-ла
на-ча-ле
на-ча-лом
на-ча-лу
на-чаль-но-го
на-чаль-ный
на-чать
на-чи-на-ет-ся
на-чи-на-ю-щи-е-ся
на-чи-на-ю-щу-ю-ся
на-ша
не-ак-тив-ная
не-без-опас-ные
не-бла-го-при-ят-ные
не-боль-шой
не-вер-ное
не-вер-ным
не-ви-ди-мых
не-го
не-дав-ние
не-дель
не-дель-ных
не-до-ста-точ-но
не-до-ступ-но
не-за-вер-шён-ная
не-за-вер-шён-ной
не-за-вер-шён-ные
не-за-вер-шён-ным
не-за-вер-шён-ны-ми
не-за-вер-шён-ных
не-за-ви-си-мая
не-за-ви-си-мо
не-за-ви-си-мо-го
не-за-ви-си-мое
не-за-ви-си-мой
не-за-ви-си-мом
не-за-ви-си-мую
не-за-ви-си-мые
не-за-ви-си-мым
не-за-ви-си-мы-ми
не-за-ви-си-мых
не-за-пу-щен-ны-ми
не-из-вест-ный
не-из-вест-ным
не-из-мен-ные
не-из-мен-ным
не-из-мен-ны-ми
не-из-мен-ных
не-из-менён-ные
ней-ро-се-те-вой
ней-ро-сеть
не-на-блю-дав-ших-ся
не-на-блю-да-е-мо-го
не-надёж-ную
не-нор-ми-ро-ван-ной
не-нуж-ные
не-ну-ле-вой
не-об-хо-ди-мо
не-об-хо-ди-мое
не-об-хо-ди-мой
не-об-хо-ди-мо-сти
не-об-хо-ди-мую
не-об-хо-ди-мых
не-огра-ни-чен-ная
не-од-но-род-ные
не-опре-делён-но-го
не-опре-делён-но-сти
не-опре-делён-ность
не-под-дер-жи-ва-е-мые
не-пол-ных
не-по-сред-ствен-но
не-прав-до-по-доб-ный
не-сколь-ки-ми
не-сколь-ких
не-сколь-ко
не-сут
не-точ-ным
не-устой-чив
не-устой-чи-вой
не-устой-чи-вость
ни-же
ниж-ней
низ-кую
но-вая
но-вой
но-вые
но-вый
но-вы-ми
но-вых
ноль
но-мер
но-ми-наль-ном
нор-маль-но-го
нор-маль-ной
нор-ма-тив-ные
нор-ми-ро-ва-на
нор-ми-ро-ван-ная
нор-ми-ро-ван-ной
нор-ми-ро-ван-ную
нор-ми-ро-ван-ные
нор-ми-ру-ют-ся
нуж-но
нуж-ны
ну-ле-вом
ну-лю
ну-ля
нулём
ны-ми
обе-их
обес-пе-чи-ва-ет
обес-пе-чи-ва-ют
об-ла-сти
об-на-ру-же-ние
об-на-ру-же-ния
об-на-ру-жен-ная
об-на-ру-жен-ной
об-на-ру-жен-ную
об-на-ру-жен-ные
об-на-ру-жен-ных
обо-зна-ча-е-мое
обо-зна-ча-ют
обо-зна-че-на
обо-зна-чен-ное
обо-их
об-ос-но-ва-ние
об-ос-но-ван-ные
об-ра-бот-ка
об-ра-бот-ке
об-ра-бот-ки
об-ра-бот-ку
об-ра-зом
об-ра-зу-ет
об-ра-зу-ют
обрат-ная
обрат-но
обрат-ное
обрат-ной
обрат-ные
обрат-ных
об-слу-жи-ва-ние
об-слу-жи-ва-нию
об-слу-жи-ва-ния
об-су-жде-ние
обу-ча-е-мые
обу-ча-ет
обу-ча-ет-ся
обу-ча-ют
обу-ча-ю-щая
обу-ча-ю-ще-го
обу-ча-ю-щей
обу-ча-ю-щие
обу-ча-ю-щий
обу-ча-ю-щим
обу-ча-ю-щих
обу-че-ние
обу-че-нии
обу-че-ний
обу-че-нию
обу-че-ния
обу-че-ни-ям
обу-че-ни-ях
обу-чен-ная
обу-чен-но-го
обу-чен-ное
обу-чен-ной
обу-чен-ную
обу-чен-ные
обу-чен-ный
обу-чен-ным
обу-чен-ны-ми
обу-чен-ных
об-щая
об-ще-го
об-щее
об-щей
об-щие
об-щий
об-щим
об-щих
об-щую
объ-еди-не-ние
объ-еди-ня-ет
объ-еди-ня-ют
объ-еди-ня-ют-ся
объ-ект
объ-яс-ня-ет
объ-ём
объ-ёма
обыч-но-го
обыч-ной
обя-за-тель-но
обя-за-тель-ные
огра-ни-че-ние
огра-ни-че-ни-ем
огра-ни-че-ний
огра-ни-че-ния
огра-ни-че-ни-я-ми
огра-ни-чен-ную
огра-ни-чен-ных
огра-ни-че-ны
огра-ни-чи-ва-ет
огра-ни-чи-ва-ет-ся
огра-ни-чи-ва-ют
один
оди-на-ко-вое
оди-на-ко-вую
оди-на-ко-вы
оди-на-ко-вые
оди-на-ко-вый
оди-на-ко-вым
оди-на-ко-вых
оди-ноч-ные
од-на
од-на-ко
од-ним
од-но
од-но-вре-мен-но
од-но-вре-мен-ных
од-но-го
од-но-знач-но
од-ной
од-но-му
од-ну
ожи-да-е-мо-го
ожи-да-е-мой
ожи-да-е-мые
ожи-да-ния
ожи-да-ют
озна-ча-ет
округ-ле-ния
округ-лён-ная
округ-лён-ный
опас-ных
опер
опе-ра-ции
опе-ра-ций
опе-ра-ци-он-ный
опе-ра-цию
опе-ра-ция
опе-ра-ци-ям
опе-ра-ци-я-ми
опе-ра-ци-ях
опе-ре-жать
опе-ре-же-ние
опе-ре-же-ни-ем
опи-са-ние
опи-сан-ные
опи-са-тель-ная
опи-са-тель-ное
опи-са-тель-ной
опи-са-тель-ную
опи-са-тель-ные
опи-са-тель-ный
опи-са-тель-ным
опи-са-тель-ны-ми
опи-са-тель-ных
опи-сы-ва-ют
опор-ную
опор-ных
опре-де-ле-нию
опре-де-ле-ния
опре-де-ля-ет
опре-де-ля-ет-ся
опре-де-ля-ют
опре-делён-ное
опре-делён-ной
опре-делён-ны-ми
оп-ти-ма-лен
оп-ти-маль-но-го
оп-ти-маль-ный
оп-ти-маль-ным
оп-ти-ми-за-ци-ей
оп-ти-ми-за-ции
оп-ти-ми-зи-ру-ют
опу-сто-ша-ет-ся
опу-сто-ше-ние
опыт
опы-тах
ор-га-ни-за-ции
ор-га-ни-за-ций
ор-га-ни-за-ци-он-ная
ор-га-ни-за-ци-он-ны-ми
ор-га-ни-за-ци-ях
ори-ги-на-лу
ори-ен-ти-ро-ван-ной
осно-ва-ны
осно-ве
основ-ная
основ-но-го
основ-ной
основ-ном
основ-ную
основ-ные
основ-ным
основ-ны-ми
основ-ных
осно-ву
оста-ва-ясь
оста-вить
остав-ля-ет
остав-ля-ют
остав-ши-е-ся
остав-ший-ся
осталь-ные
осталь-ных
оста-нов-ка-ми
остат-ка
остат-ков
оста-ток
оста-точ-ную
оста-ют-ся
остаёт-ся
осу-ще-стви-мость
от-бор
от-бо-ра
от-бо-ром
от-бо-ру
от-ве-ча-ет
от-ве-ча-ют
от-даёт
от-дель-ная
от-дель-но
от-дель-но-го
от-дель-ное
от-дель-ной
от-дель-ные
от-дель-ный
от-дель-ных
от-де-ля-ет
от-де-ля-ют
от-делён
от-ка-за
от-ка-зах
от-ка-зы
от-кло-не-ние
от-кло-не-нию
от-кло-ня-ет
от-кло-ня-ют-ся
от-клю-че-ние
от-клю-че-нии
от-кры-тые
от-кры-тых
от-ли-ча-ет-ся
от-ли-ча-ют-ся
от-ло-жен-ных
от-ме-ча-ют
от-но-си-тель-но
от-но-си-тель-ное
от-но-сит-ся
от-но-сят-ся
от-но-ся-щих-ся
от-но-ше-ние
от-но-ше-ний
отоб-ра-же-ния
отож-де-ствляя
от-ра-жа-ют
от-ри-ца-тель-ная
от-ри-ца-тель-ное
от-ри-ца-тель-ные
от-ри-ца-тель-ный
от-сут-ствие
от-сут-стви-ем
от-сут-ствии
от-сут-ствия
от-сут-ству-ет
от-счёта
от-счётам
от-счёта-ми
от-счёты
офлайн
охва-та
охва-ты-ва-ет
охва-ты-ва-ют
охва-чен-ных
оце-ни-ва-ем
оце-ни-ва-ет
оце-ни-ва-ет-ся
оце-ни-вать
оце-ни-ва-ют
оце-ни-ва-ют-ся
оце-нить
оцен-ка
оцен-ке
оцен-ки
оцен-кой
оцен-ку
оце-нок
оценён-ное
оценён-ной
оценён-но-му
оценён-ную
оценён-ных
оче-ре-дей
оче-ре-ди
оче-редь
оче-ре-дях
ошиб-ка
ошиб-кам
ошиб-ке
ошиб-ки
ошиб-кой
ошиб-ку
оши-бок
оши-боч-ную
оши-боч-ный
па-ралл
па-рал-лель-но-го
па-рал-лель-ное
па-рал-лель-ной
па-рал-лель-но-сти
па-рал-лель-ную
па-рал-лель-ные
па-рал-лель-ным
па-рал-лель-ны-ми
па-рал-лель-ных
па-ра-мет-ра-ми
па-ра-мет-ров
па-ра-мет-ром
па-ра-мет-ры
пар-ная
пар-ной
пар-ность
пар-ные
пар-ный
пар-ных
па-ру
па-ры
пер-вой
пер-вые
пе-ре-бор
пе-ре-вод
пе-ре-во-де
пе-ре-вы-бор-кам
пе-ре-вы-бо-рок
пе-ре-гру-жен-ной
пе-ре-груз-ка
пе-ре-груз-ке
пе-ре-груз-ки
пе-ре-груз-ку
пе-ре-гру-зок
пе-ред
пе-ре-мен-ной
пе-ре-ме-ща-ет-ся
пе-ре-ме-ще-ни-ем
пе-ре-на-стра-и-ва-ют
пе-ре-нос
пе-ре-но-са
пе-ре-но-сит
пе-ре-обоб-ще-ния
пе-ре-се-ка-ю-щие
пе-ре-се-ка-ю-щи-ми-ся
пе-ре-ход
пе-ре-хо-да-ми
пе-ре-хо-дов
пе-ре-хо-дом
пер-со-на-ла
пет-ри
пи-лот
пи-ло-та
пи-ло-тах
пи-ло-ты
план
пла-ни-ро-ва-ние
пла-ни-ро-ва-ни-ем
пла-ни-ро-ва-нию
пла-ни-ро-ва-ния
пла-ни-ро-воч-ные
пла-ни-ро-воч-ны-ми
пла-ни-ро-воч-ных
пла-ни-ров-щик
пла-ни-ров-щи-ка
пла-ни-ров-щи-ки
пла-ни-ров-щи-ком
пло-хо
по-ве-де-ние
по-ве-де-нии
по-ве-де-ния
по-ве-ден-че-ских
по-втор
по-вто-ра
по-вто-рам
по-вто-ра-ми
по-вто-ре
по-вто-ре-ние
по-вто-ре-ний
по-вто-ре-ния
по-вто-ре-ни-ям
по-втор-ная
по-втор-но
по-втор-но-го
по-втор-ное
по-втор-ной
по-втор-ную
по-втор-ные
по-вто-ров
по-вто-ры
по-вто-ря-е-мые
по-вто-ря-ет
по-вто-ря-ют
по-вы-ша-ет
по-гра-нич-ным
по-дав-ля-ют
под-го-тов-ка
под-го-тов-ке
под-го-тов-ки
под-го-тов-ку
под-го-тов-лен-ная
под-го-тов-лен-ное
под-го-тов-лен-ной
под-го-тов-лен-ную
под-го-тов-лен-ные
под-дер-жа-ние
под-дер-жи-ва-е-мый
под-дер-жи-ва-ет
под-дер-жи-вать
под-дер-жи-ва-ют
под-дер-жи-ва-ют-ся
под-класс
под-креп-ле-ни-ем
под-ме-на
под-про-цес-сы
под-счёт
под-твер-жда-ет
под-твер-жда-ют
под-твер-жда-ю-щие
под-твер-жда-ю-щий
под-твер-жда-ю-щи-ми
под-твер-жда-ю-щих
под-твер-жде-но
под-твер-ждён-ным
под-хо-ды
поз-во-ля-ет
позд-ней
по-зи-ци-я-ми
по-иск
по-ис-ка
по-ис-ке
по-ка-за-на
по-ка-за-но
по-ка-за-ны
по-ка-за-те-ли
по-ка-за-тель
по-ка-за-те-ля-ми
по-ка-зы-ва-ет
по-ка-зы-ва-ет-ся
по-ка-зы-ва-ют
по-кры-тие
по-кры-тия
по-лез-на
по-лез-но
по-лез-ной
по-лез-ную
по-лез-ный
по-ли-тик
по-ли-ти-ка
по-ли-ти-кам
по-ли-ти-ках
по-ли-ти-ки
по-ли-ти-ку
пол-но-го
пол-ное
пол-ной
пол-но-му
пол-но-стью
пол-но-та
пол-но-ты
пол-ную
пол-ные
пол-ный
пол-ны-ми
по-ло-вин-ную
по-ло-ви-ны
по-ло-же-ние
по-ло-жи-тель-но-го
по-ло-жи-тель-ные
по-ло-жи-тель-ный
по-ло-са
по-лу-ча-ет
по-лу-ча-ет-ся
по-лу-ча-ют
по-лу-че-ния
по-лу-че-ны
поль-за
поль-зу
поль-зы
по-ля
по-ми-мо
по-мо-гать
по-мочь
по-мо-щью
по-ни-ма-ет-ся
по-ни-ма-ния
по-прав-ки
по-прав-кой
по-прав-ку
по-ров-ну
по-рог
по-ро-га
по-ро-гах
по-ро-ге
по-ро-ги
по-ро-гов
по-ро-го-вые
по-ро-гом
по-ро-дить
по-ро-жда-ю-щих
по-ряд-ка
по-ряд-ке
по-ряд-ком
по-ряд-ку
по-ря-док
по-сколь-ку
по-сле
по-след
по-след-нее
по-след-ние
по-след-ний
по-сле-до-ва-тель-ная
по-сле-до-ва-тель-но-го
по-сле-до-ва-тель-ное
по-сле-до-ва-тель-ной
по-сле-до-ва-тель-ном
по-сле-до-ва-тель-но-стей
по-сле-до-ва-тель-но-сти
по-сле-до-ва-тель-ность
по-сле-до-ва-тель-ную
по-сле-до-ва-тель-ные
по-сле-до-ва-тель-ный
по-сле-до-ва-тель-ных
по-след-ствие
по-след-ствия
по-сле-ду-ю-щая
по-сле-ду-ю-ще-го
по-сле-ду-ю-щей
по-сле-ду-ю-щие
по-сле-ду-ю-щим
пост
по-ста-нов-ке
по-сто-ян-ная
по-сто-ян-но
по-сто-ян-но-го
по-сто-ян-ное
по-сто-ян-ном
по-сто-ян-ным
по-стро-е-нии
по-ступ-ле-ние
по-ступ-ле-нии
по-ступ-ле-ний
по-ступ-ле-ния
по-тен-ци-аль-ные
по-те-ри
по-терь
по-ток
по-то-ке
по-то-ки
по-тре-бо-вать-ся
по-че-му
по-чти
по-это-му
по-яв-ля-ют-ся
прав-до-по-до-бие
пра-вил
пра-ви-ла
пра-ви-ла-ми
пра-ви-ло
пра-ви-лом
пра-ви-лу
пра-виль-но
пра-виль-ный
пре-вос-хо-дит
пре-вос-ход-ства
пре-вос-ход-ство
пре-вы-ша-ет
пре-вы-ше-ние
пре-вы-ше-ний
пре-вы-ше-ния
пре-де-ла-ми
пре-де-лы
пре-дик-тив-ная
пре-дик-тив-но-го
пре-дик-тив-ной
пре-дик-тив-ную
пре-дик-тив-ные
пре-дик-тив-ный
пре-дик-тив-ным
пред-кри-зис-но-го
пред-кри-зис-но-му
пред-кри-зис-ный
пред-ла-га-ют
пред-ло-жен-но-го
предот-вра-ща-ет
предот-вра-ще-ние
предот-вра-ще-ний
предот-вра-ще-нию
предот-вра-ще-ния
предот-вра-ще-но
пред-пи-са-ние
пред-пи-сы-ва-ю-ще-го
пред-пи-сы-ва-ю-щий
пред-по-ла-га-е-мом
пред-по-ла-га-е-мые
пред-по-ла-га-е-мый
пред-по-ла-га-ет
пред-по-ла-гая
пред-по-ло-же-ний
пред-по-ло-же-ния
пред-по-ло-же-ни-я-ми
пред-по-сыл-ка-ми
пред-по-чи-та-ет
пред-по-чте-ния
пред-ска-зан-ных
пред-ска-зы-ва-ет
пред-ста-вить
пред-став-ле-ние
пред-став-ле-нии
пред-став-ле-ния
пре-ду-пре-жде-ния
пред-ше-ство-ва-ли
пред-ше-ство-ва-ло
пред-ше-ству-ет
пред-ше-ству-ю-щая
пред-ше-ству-ю-щих
пре-ды-ду-щей
пре-ды-ду-щем
пре-ды-ду-щие
преж-ней
преж-нем
преж-ние
преж-ний
преж-них
преж-няя
пре-иму-ще-ства
пре-иму-ще-ство
пре-иму-ще-ству
пре-кра-ща-ют-ся
пре-об-ра-зо-ва-ние
пре-об-ра-зо-ва-ния
пре-фик-сов
при-бли-жа-ет
при-бли-же-ние
при-бли-же-ния
при-ве-де-ны
при-во-дит
при-во-дит-ся
при-вя-зан-ная
при-год-но-сти
при-ем-ле-мое
при-знак
при-зна-ка
при-зна-кам
при-зна-ки
при-зна-ков
при-ло-же-ние
при-ло-же-нии
при-ме-не-ние
при-ме-не-нии
при-ме-не-ния
при-ме-ня-ет
при-мер
при-мер-но
при-ни-ма-ем
при-ни-ма-ет-ся
при-ня-тия
при-ня-тые
при-о-ри-тет
при-пи-сы-вая
при-спо-соб-ле-ние
при-хо-дит-ся
при-чин-но-го
при-чин-ные
при-чин-ный
причём
про-бле-ма
про-бле-му
про-бы
про-ве-рен-ной
про-ве-рен-ном
про-ве-рен-ную
про-ве-рен-ный
про-ве-рен-ных
про-ве-рить
про-вер-ка
про-вер-ках
про-вер-ке
про-вер-ки
про-вер-ку
про-ве-рок
про-ве-ря-ем
про-ве-ря-е-мое
про-ве-ря-ет
про-ве-ря-ет-ся
про-ве-ря-ют
про-ве-ря-ют-ся
про-ве-сти
про-гноз
про-гно-за
про-гно-зи-ро-ва-ние
про-гно-зи-ру-ем
про-гно-зи-ру-ет
про-гноз-ная
про-гноз-ной
про-гно-зом
про-гно-зу
про-го-нах
про-го-нов
про-длён-ной
про-длён-ным
про-дол-же-ни-ем
про-дол-же-ния
про-ек-ти-ро-ва-ни-ем
про-игрыш
про-из-во-ди-тель-но-сти
про-изо-шла
про-ис-хо-дит
про-пус-ки
про-пу-щен-ные
про-смот-рен-ные
про-стое
про-стой
про-стые
про-стым
про-сты-ми
про-тив
про-то-кол
про-то-ко-ла
про-то-ко-лы
про-хо-да
про-хо-дит
про-хо-дов
про-хо-дят
про-хо-дя-щая
про-хо-жде-ние
про-це-ду-ра
про-це-ду-ру
про-це-ду-ры
про-цент
про-цен-тах
про-цен-тиль
про-цен-тиль-ные
про-цен-ти-ля-ми
про-цен-ты
про-цесс
про-цес-са
про-цес-сам
про-цес-сах
про-цес-се
про-цесс-но
про-цесс-ной
про-цес-сов
про-цес-су
про-шли
про-яви-лась
пря-мое
пуас-со-нов-ские
пу-лам
пу-ла-ми
пу-лов
пу-лом
пу-лы
пу-ти
пять
ра-бо-та
ра-бо-та-ю-щей
ра-бот-ник
ра-бот-ни-ка
ра-бот-ни-ков
ра-бо-ту
ра-бо-ты
ра-бо-че-го
ра-бо-чем
ра-бо-чие
ра-бо-чий
ра-вен
ра-вен-ства
ра-вен-ство
рав-на
рав-но
рав-ной
рав-ны
рав-ные
рав-ны-ми
раз-брос
раз-ветв-ле-ния
раз-вил-ки
раз-гру-жа-ет-ся
раз-дел
раз-де-ле-ние
раз-де-ле-ны
раз-де-лы
раз-дель-но
раз-дель-ные
раз-де-ля-ет
раз-де-ля-ют
раз-ли-ча-ет
раз-ли-чать
раз-ли-ча-ют
раз-ли-ча-ют-ся
раз-ли-ча-ю-щим-ся
раз-ли-че-ния
раз-ли-чие
раз-ли-чий
раз-ли-чи-мых
раз-ли-чия
раз-ли-чи-я-ми
раз-лич-ных
раз-ло-же-ние
раз-ме-ща-ли
раз-но-стей
раз-но-сти
раз-ность
раз-но-стью
раз-ные
раз-ных
раз-ра-бот-ки
раз-ре-ша-ют-ся
ран-жи-ро-ва-ние
ран-жи-ро-ва-нию
ран-жи-ро-ва-ния
ран-ние
рань-ше
рас-пи-са-ние
рас-пи-са-ния
рас-пре-де-ле-ние
рас-пре-де-ле-ния
рас-про-странён-ность
рас-смат-ри-ва-ем
рас-смат-ри-ва-е-мой
рас-смат-ри-ва-ет
рас-смат-ри-ва-ют-ся
рас-смот-ре-ние
рас-смот-рен-ные
растёт
рас-хо-дя-щи-е-ся
рас-хо-жде-ние
рас-ши-ре-ние
ре-а-ги-ро-ва-ния
ре-ак-ци-ей
ре-ак-ции
ре-ак-ций
ре-ак-цию
ре-ак-ция
ре-ак-ци-ях
ре-а-ли-за-ции
ре-а-ли-за-ция
ре-а-ли-зо-ван-ная
ре-а-ли-зо-ван-ное
ре-а-ли-зо-ван-ную
ре-а-ли-зо-ван-ные
ре-а-ли-зо-ван-ным
ре-а-ли-зо-ван-ных
ре-аль-ное
ре-аль-ной
ре-аль-ные
ре-аль-ный
ре-аль-ных
ре-ги-стри-ру-ет
ре-ги-стри-ру-ют-ся
ре-грес-сия
ред-ких
ре-жи-ма
ре-зер-ви-ру-ют-ся
ре-зуль-тат
ре-зуль-та-та
ре-зуль-та-там
ре-зуль-та-тах
ре-зуль-та-тов
ре-зуль-та-ты
ре-ко-мен-да-ция
ре-сурс
ре-сур-са
ре-сур-сам
ре-сур-са-ми
ре-сур-сах
ре-сурс-ная
ре-сурс-ные
ре-сурс-ный
ре-сурс-ны-ми
ре-сурс-ных
ре-сур-сов
ре-сур-сы
ре-тро-спек-тив-ная
ре-тро-спек-тив-но
ре-тро-спек-тив-ной
ре-ша-ет
ре-ша-ют
ре-ше-ние
ре-ше-нии
ре-ше-ний
ре-ше-нию
ре-ше-ния
ре-ше-ни-я-ми
риск
рис-ка
ри-су-нок
рост
ро-стом
ру-ко-пи-си
ру-ко-пись
рус-ский
са-ма
са-ми
са-мо-под-дер-жи-ва-ю-щий-ся
са-мо-про-вер-ке
сат-тер-ту-эй-та
све-де-ния
све-де-ни-ях
сви-де-тельств
сви-де-тель-ства
сви-де-тель-ства-ми
сви-де-тель-ство
сво-бод-ным
сво-бо-ды
свое-вре-мен-ность
свой-ства
свя-зан-ные
свя-за-но
свя-зей
свя-зи
свя-зы-ва-ет
свя-зы-ва-ют
связь
свя-зью
свя-зях
сге-не-ри-ро-ван-ных
се-лек-тор
се-лек-то-ра
се-ман-ти-ка
се-мей-ства
се-мей-стве
се-мей-ство
се-рии
се-рия
се-рые
се-ти
сет-ка
сет-ке
сет-ку
се-тью
сиг-нал
сиг-на-ла
сиг-на-лов
си-ле
си-лы
силь-но-го
силь-ное
силь-ные
силь-ный
силь-ны-ми
си-му-ли-ру-ет
си-му-ля-тор
си-му-ля-то-ра-ми
си-му-ля-то-ров
си-му-ля-ци-ей
си-му-ля-ции
си-му-ля-ций
си-му-ля-ци-он-ные
си-му-ля-ция
си-му-ля-ци-ям
си-му-ля-ци-ях
син-те-ти-че-ские
син-те-ти-че-ский
син-те-ти-че-ских
син-хро-ни-за-ци-ей
син-хро-ни-за-ции
син-хро-ни-за-ция
си-стем
си-сте-ма
си-сте-ма-ми
си-сте-ме
си-стем-ное
си-стем-ные
си-сте-му
си-сте-мы
сквоз-ную
скон-стру-и-ро-ван-ные
скон-стру-и-ро-ван-ных
скор-рек-ти-ро-ван-ный
скор-рек-ти-ро-ван-ных
скрип-та-ми
скрип-ты
сла-бо
сла-га-е-мое
след
сле-до-ва-тель-но
сле-до-вать
сле-ду-ет
сле-ду-ют
сле-ду-ю-ще-го
сле-ду-ю-щей
сле-ду-ю-щие
сле-ду-ю-щий
сле-ду-ю-щих
сле-ду-ю-щую
сло-ва
слож-ной
слож-но-сти
слож-ность
слож-ные
слож-ных
сло-я-ми
слу-жат
слу-жит
слу-чае
слу-ча-ем
слу-чай
слу-чай-ность
слу-чай-но-стью
слу-чай-ные
слу-ча-ях
смеж-ные
сме-ну
сме-си
смысл
сна-ча-ла
сни-жа-ет
сни-жа-ет-ся
сни-же-ние
сни-же-ния
сни-зить
со-ав-то-ры
со-блю-да-ют-ся
со-бы-тие
со-бы-ти-ем
со-бы-тий
со-бы-тию
со-бы-тия
сов-мест-но-го
сов-мест-ное
сов-мест-ном
сов-мест-ный
сов-па-да-ет
сов-па-да-ют
сов-па-да-ю-щие
сов-па-де-ние
со-гла-со-ва-ния
со-гла-со-ван-ный
со-гла-со-ван-ны-ми
со-гла-со-ва-ны
со-дер-жа-нию
со-дер-жат
со-дер-жит
со-зда-ние
со-кра-ще-ние
со-кра-ще-ни-ем
со-кра-ще-нии
со-кра-ще-нию
со-кра-ще-ния
со-об-ща-ет-ся
со-от-вет-ствен-но
со-от-вет-стви-ем
со-от-вет-ствия
со-от-вет-ству-ют
со-от-вет-ству-ю-щая
со-от-вет-ству-ю-щее
со-от-вет-ству-ю-щей
со-от-вет-ству-ю-щие
со-от-вет-ству-ю-щий
со-от-вет-ству-ю-щих
со-от-вет-ству-ю-щую
со-по-став-лен-ном
со-по-став-ля-ют
со-по-став-ля-ют-ся
со-про-во-ждать-ся
со-став
со-ста-ва
со-ста-ви-ла
со-став-ля-ет
со-став-ля-ют
со-сто-ит
со-сто-я-ние
со-сто-я-ний
со-сто-я-нию
со-сто-я-ния
со-сто-я-ни-ям
со-сто-я-ни-ях
со-хра-не-на
со-хра-не-ние
со-хра-не-нии
со-хра-не-ния
со-хра-не-ны
со-хра-нив
со-хра-нить
со-хра-ня-ем
со-хра-ня-ет
со-хра-ня-ет-ся
со-хра-нять
со-хра-ня-ют
со-хра-ня-ют-ся
со-хра-ня-ю-щей
со-хра-няя
со-хранён
со-хранён-ное
со-хранён-ной
со-хранён-ную
со-хранён-ные
со-хранён-ным
со-хранён-ных
со-че-та-ет
со-че-та-ет-ся
со-че-та-ние
со-че-та-ний
спе-ци-фи-ка-ции
спо-кой-ная
спо-кой-ном
спо-кой-ных
спо-со-бен
спо-соб-на
спо-соб-но
сра-ба-ты-ва-ет
сра-ба-ты-ва-ние
сра-ба-ты-ва-ний
сра-ба-ты-ва-ния
сра-бо-тал
срав-не-ние
срав-не-нии
срав-не-ний
срав-не-ния
срав-не-ни-я-ми
срав-не-ни-ях
срав-ни-ва-ем
срав-ни-ва-ет
срав-нить
сре-ди
сред-нее
сред-ней
сред-не-квад-ра-тич-но-му
сред-ние
сред-ний
сред-ним
сред-них
сред-нюю
сред-няя
сре-ды
ста-дии
стан-дарт-но-го
стан-дарт-ное
стан-дарт-ном
стан-дарт-ный
ста-но-вит-ся
старт
стар-та
стар-те
стар-том
стар-ты
стат
ста-ти-сти-ки
ста-ти-сти-че-ских
ста-ти-сти-че-ско-го
ста-тич
ста-ти-че-ская
ста-ти-че-ские
ста-ти-че-ский
ста-ти-че-ским
ста-ти-че-ски-ми
ста-ти-че-ско-го
ста-ти-че-ское
ста-ти-че-ской
ста-тьи
ста-тью
ста-ци-о-нар-ный
сте-пе-ня-ми
сто-и-мо-сти
сто-и-мость
сто-и-мо-стью
сто-ит
стол-бец
сто-ят
стра-те-гий
стрес-сах
стрес-се
стрес-со-вые
стрес-сы
стро-го
стро-го-го
стро-ит
строк
стро-ка
стро-ке
стро-ят-ся
струк-тур
струк-ту-ра
струк-ту-рам
струк-ту-ре
струк-тур-ная
струк-тур-но-го
струк-тур-ное
струк-тур-ную
струк-тур-ные
струк-тур-ных
струк-ту-рой
струк-ту-ру
струк-ту-ры
стью-ден-та
сум-мар-ной
сум-ме
сум-ми-ро-ва-ния
сум-мы
суф-фик-са
су-ще-ствен-ны
су-ще-ствен-ны-ми
су-ще-ствен-ных
су-ще-ству-ю-ще-му
су-ще-ству-ю-щим
сце-на-ри-ев
сце-на-рии
сце-на-рий
сце-на-рия
сце-на-ри-ям
сце-на-ри-ях
счи-та-ет
счи-та-ет-ся
счи-тая
счёт-чи-ка-ми
счёт-чи-ки
табл
та-блиц
та-бли-ца
та-бли-цы
так-же
та-кие
та-ким
та-ких
та-ко-го
та-кой
тек-ста
тек-сте
те-ку-щие
тер-ми-наль-но-му
тер-ми-наль-ный
тест
те-ста
те-сте
те-сти-ро-ва-ние
те-сти-ро-ва-нии
те-сти-ро-ва-ния
те-сто-вая
те-сто-вой
те-сто-вом
те-сто-вые
те-сто-вым
те-сто-вы-ми
те-сто-вых
те-сту
те-сты
тех-ни-че-ская
тех-ни-че-ские
тех-ни-че-ский
те-че-ние
ти-па
ти-пов
ти-пы
то-гда
то-ке-ном
то-ке-ны
толь-ко
то-по-ло-гии
то-по-ло-гий
то-по-ло-гия
то-чек
точ-ках
точ-ки
точ-ку
точ-ная
точ-но
точ-но-го
точ-ное
точ-ной
точ-но-сти
точ-ность
точ-но-стью
точ-ным
тра-ек-то-рии
тра-ек-то-рий
тра-ек-то-рия
тра-ек-то-ри-ях
тра-пе-ций
трасс
трас-сам
трас-сы
тре-бо-ва-ния
тре-бу-ет
тре-бу-ют
тре-бу-ют-ся
тре-во-ги
тру-до-за-трат
трёх
ты-ся-чах
ты-ся-чи
тяжёлая
тяжёло-го
тяжёлое
тяжёлой
тяжёлую
тяжёлым
тяжёлых
уби-ра-ет
уве-ли-че-ние
уве-ли-че-нии
уве-ли-че-нию
уве-ли-чен-ной
уве-ли-чи-ва-ет
уве-ли-чи-ва-ет-ся
уве-ли-чи-вать
уве-ли-чи-ва-ют
уве-ли-чи-вая
уда-ле-ние
уда-ля-ют-ся
удво-ен-ной
удо-вле-тво-ря-ет
уз-кие
уз-ко-го
уз-ком
уз-ла
ука-зан-ном
ука-за-ны
улуч-ша-ет
улуч-шать
улуч-ше-ние
улуч-ше-нию
улуч-ше-ния
умень-ша-ет
умень-ша-ет-ся
умень-ша-ют
умень-ше-ние
умень-шить
умно-жа-ю-щий
умно-жа-ю-щую
уни-вер-саль-но-му
урав-не-ния
урав-не-ни-ям
уро-вень
уров-не
уси-ле-ние
уси-ле-нию
уси-ли-ва-ют
уско-рен-ное
усло-ви-ем
усло-вии
усло-вий
усло-вия
усло-ви-ям
усло-ви-ях
услов-ная
услов-но
услов-ные
услов-ных
услож-не-ние
услож-не-ния
успех
усред-ня-ет
усред-ня-ет-ся
усред-ня-ют
усред-ня-ют-ся
усред-нён-ная
уста-нав-ли-ва-ют
уста-но-ви-ли
уста-но-вить
уста-нов-ле-ние
уста-нов-ле-но
уста-нов-ле-ны
устой-чи-во-сти
устой-чи-вость
устой-чи-вую
устой-чи-вые
устра-ня-ет
утвер-жда-ет
утвер-жде-ние
утвер-жде-ния
утвер-ждён-ный
уточ-ня-ют
ухуд-ша-ет-ся
ухуд-ша-ют-ся
ухуд-шив
участ-ву-ют
уча-ща-ет-ся
учи-ты-ва-ет
учи-ты-вать
учи-ты-ва-ют
учи-ты-ва-ют-ся
учёт
учёта
учётом
фа-зо-вый
фак-ти-че-ская
фак-ти-че-ские
фак-ти-че-ское
фак-ти-че-ском
фак-ти-че-скую
фи-зи-че-ской
фик-са-ции
фик-си-ро-ван-ная
фик-си-ро-ван-но-го
фик-си-ро-ван-ное
фик-си-ро-ван-ной
фик-си-ро-ван-ном
фик-си-ро-ван-но-му
фик-си-ро-ван-ную
фик-си-ро-ван-ные
фик-си-ро-ван-ный
фик-си-ро-ван-ным
фик-си-ро-ван-ных
фик-си-ру-ет
фик-си-ру-ет-ся
фик-си-ру-ют-ся
филь-тра-цию
фиш-ки
фор-ми-ру-ет-ся
функ-ци-ей
функ-ции
функ-ция
ха-рак-те-ри-зу-ет
ха-рак-те-ри-зу-ют
хе-шей
хо-тя
хра-нит
хро-но-ло-ги-че-ские
хро-но-ло-ги-че-ский
ху-же
це-ле-вая
це-ле-вой
це-лей
це-ли
цель
цен-зу-ри-ро-ва-ни-ем
цен-но-сти
цен-ность
цен-но-стью
це-ны
цик-ла
цик-ла-ми
цик-лах
цик-ле
цик-лы
ча-сти
ча-сто
ча-сто-та
ча-сто-ту
ча-сто-ты
часть
че-го
че-ло-ве-че-ское
че-ре-ду-ю-ще-е-ся
че-рез
че-ты-ре
че-тырём
че-тырёх
чис-ла
чис-ле
чис-лен-но
чис-лен-ные
чис-ло
член
чле-нов
что-бы
чув-стви-тель-но-сти
чув-стви-тель-ность
чуть
ша-га
ша-ге
ша-гом
ше-сти
шест-на-дцать
шесть
ши-ре
ши-ро-ком
шка-лах
шка-ле
шу-ма
эври-сти-ка
эври-сти-ки
эври-сти-кой
эври-сти-ку
эк-ви-ва-лен-тах
эк-ви-ва-лент-ность
эк-ви-ва-лен-ты
эк-зем-пляр
эк-зем-пля-ра
эк-зем-пля-ров
эк-зем-пля-ры
экс-пе-ри-мент
экс-пе-ри-мен-та
экс-пе-ри-мен-таль-ный
экс-пе-ри-мен-те
экс-пе-ри-мен-тов
экс-пе-ри-мен-ты
экс-плу-а-та-ци-ей
экс-плу-а-та-ции
экс-плу-а-та-ци-он-но-го
экс-тра-по-ли-ро-вать
экс-тра-по-ля-ции
экс-тра-по-ля-цию
экс-тра-по-ля-ция
эле-мен-ты
эм-пи-рич
эм-пи-ри-че-ская
эм-пи-ри-че-ский
эм-пи-ри-че-ское
эн-до-ген-ных
эн-тро-пия
эпи-зод
эпи-зо-да
эпи-зо-дах
эпи-зо-дов
эпи-зо-ды
эта-лон
эта-ло-на
этап
эта-па-ми
эта-пе
эта-пы
этим
эти-ми
этих
этой
этом
этот
эф-фект
эф-фек-та
эф-фек-там
эф-фек-те
эф-фек-тив-но-сти
эф-фек-тив-ность
эф-фек-тов
эф-фек-том
эф-фек-ту
эф-фек-ты
яв-ля-ет-ся
яв-ля-ют-ся
яв-но
яв-ное
яв-ные
яв-ным
яв-ны-ми
яв-ных
язык
язы-ков
язы-ком
яче-ек
}

\title{Дополнительный экспериментальный протокол}
\titlerunning{Дополнительные материалы}
\author{Анонимные авторы}
\authorrunning{Анонимные авторы}
\institute{Анонимная рукопись}
\maketitle
Первые разделы сохраняют запись исходного последовательного эксперимента P0 и вспомогательных исследований. Заключительный раздел содержит независимое структурное расширение P1--P6, полные сравнения и диагностику чувствительности. Два подтверждающих протокола и последующие диагностические анализы рассматриваются раздельно.

Эти материалы описывают вспомогательные симуляционные кампании, прогнозирование на реальных журналах, полную основную оценку реакций и последующие диагностические проверки. Диагностический протокол был задан после получения основных результатов.

\section{Методика: от истории работы к подготовленной реакции}
На рис.~\ref{fig:workflow} показано, какие данные используются для обучения модели, выбора реакции и её проверки. Тестовые исходы оценивают сохранённую реакцию, не меняя её выбора.
\begin{figure}[t]
\centering\includegraphics[width=\textwidth]{generated/workflow_ru.pdf}
\caption{Поток информации в методике подготовки и её контролируемой оценке. Инженеры задают ограничения процесса и предполагаемые эффекты действий. Последний этап оценивает политику, зафиксированную на предыдущем.}\label{fig:workflow}
\end{figure}
\subsection{Представление исполнения и восстановление наблюдаемого состояния}
Под задачей понимается заявка, проходящая несколько операций.

Модель $\mathcal M=(N,\mathcal R,K,\alpha,F_D,F_R)$ задаёт их порядок через $N$, группы взаимозаменяемых ресурсов $\mathcal R$, доступное число ресурсов в каждой группе $K$, назначение операций группам $\alpha$ и функции длительности обработки и вероятности возврата. Расписания внешних поступлений задаются отдельно. Симулятор хранит состояние каждой задачи, занятость ресурсов, очереди и число незавершённых задач $w(t)$, обозначаемое WIP. Это число включает задачи как в очередях, так и в обработке. При параллельных операциях подсчёт каждой задачи один раз отличается от суммирования очередей и занятых мест отдельных операций.

Готовые операции ожидают ресурс в порядке поступления. При представлении процесса сетью Петри фишки отмечают доступность операций; входные фишки резервируются при постановке операции в очередь, выходные появляются по завершении. Занятость удовлетворяет $0\leq b_r(t)\leq K_r$. Длительность обработки выбирается один раз при её начале $t_s$ только по наблюдаемому тогда состоянию:
\begin{equation}\label{eq:duration}
\log D_a=f_a(x,\ell(t_s))+\epsilon_a,\qquad \epsilon_a\sim\mathcal N(0,\sigma^2).
\end{equation}
В момент решения о возврате при завершении $t_c$:
\begin{equation}\label{eq:return}
R\sim\mathrm{Bernoulli}\!\left(\operatorname{sigmoid}(g(x,\ell(t_c)))\right).
\end{equation}
Признаки задачи $x$ доступны в момент решения. Длина очереди операции $a$ обозначена $q_a$; признаки нагрузки включают отдельные для операций $\log(1+q_a/K_{\alpha(a)})$ и $\log(1+w/6)$. Очередь перед стартом включает начинающуюся операцию; очередь при решении о возврате ещё не включает возможную возвращаемую операцию. Будущее состояние завершения не используется как признак при старте.

Для логарифма длительности обучается ridge-регрессия ($\alpha=1$), для возвратов --- логистическая регрессия ($C=10$). Статические модели сохраняют контекст и конечную мощность, но исключают признаки нагрузки. Зависимые модели сохраняют оба типа признаков. В последовательной реализации признак возврата наблюдает очередь B при завершении C и использует зафиксированный масштаб 2; длительность B использует мощность 3. При изменении числа исполнителей вмешательством оба знаменателя сохраняют значения, использованные при обучении. В сохранённой реализации зафиксированы общее стандартное отклонение остатков и определения признаков.

Дополнительное правило выбора использует улучшения валидационной функции потерь $\Delta_D$ и $\Delta_R$: признаки нагрузки для длительности сохраняются при $\Delta_D>0.01$, для возвратов --- при $\Delta_R>0.005$. Для длительности используется отрицательное логарифмическое правдоподобие логарифма времени, для возвратов --- бинарная кросс-энтропия. Такой выбор проверяет прогноз в наблюдаемом диапазоне. Эксперимент с заранее выбранной реакцией сравнивает его с безусловным учётом нагрузки, не предполагая, что выбор по точности оптимален для подготовки реакции.

\subsection{Исследование нагрузок и выбор реакции}
Инженер задаёт возможные сценарии нагрузки $\mathcal S_p$, доступные действия $\mathcal A$ и их стоимость. Сценарии могут менять интенсивность поступлений, состав ресурсов или задач; основной эксперимент проверяет изменение интенсивности поступлений. Каждый вариант реакции объединяет действие $a$ и порог $q$ наблюдаемого числа незавершённых задач. Вмешательство начинается после двух последовательных превышений при шаге наблюдения пять единиц:
\begin{equation}\label{eq:trigger}
\tau_q=\inf\{t_k:w(t_{k-1})\geq q\ \land\ w(t_k)\geq q\}.
\end{equation}
Реализованное будущее расписание поступлений недоступно правилу. Два отсчёта подавляют одиночные превышения; это явное правило очереди, а не обученный классификатор кризиса. Подготовка предшествует возмущению; включение действия может следовать за началом всплеска и при этом опережать тяжёлую перегрузку. Мы не прогнозируем момент начала ненаблюдаемого будущего всплеска поступлений.

Проверяются отсутствие реакции, добавление одного исполнителя B и уменьшение вероятности будущих возвратов вдвое. Увеличение мощности позволяет начать дополнительные готовые операции, не меняя уже выбранные длительности. Сокращение возвратов действует только на последующие решения. Для горизонта $H$:
\begin{equation}\label{eq:cost}
J(\pi)=\int_0^H w_\pi(t)\,dt+c_a(H-\tau_q)_+,
\end{equation}
причём стоимость действия равна нулю, если порог не сработал. Интеграл числа незавершённых задач измеряет время их нахождения в системе; реализация приближает его методом трапеций по отсчётам с шагом пять. Стоимости $c_{\mathrm{capacity}}=2$ и $c_{\mathrm{rework}}=0.5$ заданы в условных эквивалентах времени задачи, а не деньгах. Все планировщики получают одинаковые каталог, пороги и стоимости.

Каждая обученная модель симулирует библиотеку планирования и выбирает
\begin{equation}\label{eq:policy}
\widehat\pi=\arg\min_{\pi\in\Pi}\frac{1}{|\mathcal S_p|L}
\sum_{s\in\mathcal S_p}\sum_{l=1}^{L}J_{\widehat{\mathcal M},s,l}(\pi).
\end{equation}
Равенства разрешаются по фиксированному порядку каталога. Общие случайные потоки уменьшают случайные различия между моделями и действиями. Выбранная политика сохраняется до генерации любой тестовой траектории. Применение правила и оценка эффекта выполняются затем.

\subsection{Проверка охвата данных и результата решения}
Перед использованием подготовленной реакции проверяются наблюдаемый диапазон нагрузок, улучшение прогноза на валидационных данных, предположения о ресурсах и предполагаемые эффекты действий. Если два механизма совпадают на всех наблюдавшихся состояниях, но различаются за их пределами, одной истории недостаточно для различения будущих реакций. Поэтому экстраполяция требует предположений или дополнительных измерений. Явное моделирование нагрузки не устраняет это ограничение идентификации.

Мы сохраняем все калибровочные траектории и проверяем каждую на наличие длительной тяжёлой перегрузки. В эксперименте с заранее выбранной реакцией тяжёлое событие означает $w(t)\geq100$ в течение 200 последовательных единиц времени с точностью шага наблюдения. Запас времени сигнала измеряется относительно начала события в парной траектории без действия. Срабатывание без всплеска показывается отдельно от ложной тревоги: сама базовая работа может быть неустойчивой. Полезный сигнал не обязательно предотвращает событие, поэтому измеряются также стоимость и частота событий после действия.

Результат применения методики включает сохранённую модель, рассмотренные сценарии нагрузки, подготовленные действие и порог его включения, а также данные о проверенном рабочем диапазоне. Для применения в работающей организации также нужно проверить эффекты вмешательств и ограничения персонала, затем провести контролируемую оценку внедрения. Заданные в симуляциях эффекты действий не дают таких свидетельств.

\section{План оценки}
Метод исследования сочетает создание инструмента, контролируемые симуляции, выявление структуры и хронологический прогноз на открытых журналах. Эксперимент с заранее выбранной реакцией является основным; предыдущие серии дают вспомогательные свидетельства. Единицами неопределённости служат независимые обучающие наборы, а не повторные траектории одной обученной модели.

\subsection{Контролируемый процесс и эксперимент с заранее выбранной реакцией}
Процесс имеет вид A$\rightarrow$B$\rightarrow$C с выходом либо возвратом в B. Мощности равны $(2,3,2)$, базовые средние длительности --- $(1,2.5,1.4)$. Пуассоновские поступления несут сложность $c\sim\mathrm{Bernoulli}(0.35)$, умножающую среднее обслуживание на $1+0.35c$. Логарифмическое стандартное отклонение логнормального шума равно 0.25. Четыре мира включают отсутствие обратной связи (фиксированный), влияние нагрузки на обслуживание, влияние на возвраты либо оба механизма. Активные механизмы:
\begin{align}\label{eq:truth}
\mathbb E[D_B\mid q,c]&=2.5(1+0.35c)\left(1+\frac{1.2q}{q+6}\right),\\
P(R=1\mid w,c)&=0.04+0.03c+\frac{0.45z^2}{1+z^2},\quad z=\frac{\max(0,w-6)}{18}.
\end{align}
Неактивная связь убирает соответствующее слагаемое нагрузки. Обучаемые логлинейная и логистическая функции отличаются от порождающих зависимостей. Симуляции моделей ограничивают логарифм длительности диапазоном $[-4,5]$, вероятность возврата --- $[0.001,0.95]$; заданная истинная система не заменяется обученным прогнозом.

Новый эксперимент E7 использует 12 независимо сгенерированных обучающих наборов в каждом мире. Восемь калибровочных эпизодов дважды используют интенсивности $\{0.25,0.40,0.55,0.65\}$; три валидационных --- $\{0.30,0.50,0.60\}$. Поступления заканчиваются в момент 300, затем каждый эпизод полностью опустошается. Все эпизоды сохраняются в оценке; утверждение о докризисной калибровке требует отсутствия тяжёлых событий в аудите рассматриваемой группы.

Библиотека планирования сочетает базовые интенсивности 0.40/0.50 с отсутствием всплеска, интенсивностью 1.10 на $[200,300)$ либо 1.50 на $[200,400)$. Горизонт равен 1,400, число повторов --- два. Кандидатные пороги: 20, 40 и 80. Каждая модель выбирает одну политику по шести условиям. В тесте используются новые базы 0.45/0.55 и отсутствие всплеска, интенсивность 1.35 на $[350,500)$ либо 1.65 на $[350,600)$, горизонт 1,800 и восемь повторов. Таким образом, момент и величина всплеска отличаются от подготовки. Случайные потоки калибровки, планирования и теста разделены. Случайность операций, привязанная к их идентификаторам, образует парные действия независимо от порядка событий.

Основной кандидат безусловно учитывает нагрузку. Сохраняются статическое планирование и модель с валидационным выбором, а также отсутствие действия, эвристика добавления мощности при WIP 20 и эвристика сокращения возвратов при WIP 20. Последние две строго проверяют, даёт ли обучение выбору действия дополнительную пользу. Пилот разработки использовал одно отдельное зерно на мир и исключён из статистических выводов. Его рассмотрение предшествовало фиксации протокола; оно мотивировало основной кандидат и общую случайность планирования. Основная серия включает 48 обученных наборов, 12,096 планировочных и 16,128 тестовых симуляций. Тестовые исходы основной серии не меняют выбор моделей или протокол.

Основной показатель --- стоимость, усреднённая с равными весами по шести условиям и восьми повторам. Основным механизмом являются возвраты, выбранные на основе предшествующих экспериментов. Два главных сравнения сопоставляют планирование с нагрузкой со статическим планированием и добавлением мощности при 20; используются парные различия по независимым обучающим наборам и двусторонние 97.5\%-ные интервалы Стьюдента с 11 степенями свободы, что даёт поправку Бонферрони на два утверждения. Сокращение возвратов при 20 всегда показано, включая равенство результатов. Остальные 95\%-ные интервалы описательные. Дополнительное сравнение усредняет внутри каждого зерна стоимость каждого мира, нормированную на отсутствие действия, с равными весами механизмов. Анализ чувствительности использует множители стоимости 0.5, 1 и 2; политика заново выбирается только по сохранённым результатам планирования, затем оценивается на сохранённых тестовых действиях.

\subsection{Вспомогательные проверки симуляций и реальных журналов}
Предыдущие модели с заданной структурой используют шесть независимых прогонов на мир, обучающие интенсивности 0.40, 0.65, 0.80, 1.05, 1.30 и валидационные 0.50, 0.90, 1.15, поступления до момента 400 и полное опустошение. Начальный аудит E1 проверяет 235,114 операций и 432 симуляции. E5 повторно использует эти модели с независимой внутренней случайностью, базами 0.40/0.50/0.65, возможным всплеском 1.50 на $[300,500)$, горизонтом 2,000 и восемью повторами. Сильный эмпирический статический базовый метод повторно выбирает логарифмы длительностей внутри классов операции и сложности. Нейросетевой вариант с двумя слоями по 32 tanh-узла получает тот же контекст и нагрузку. E5 также проверяет B+1 и половинную вероятность возвратов при базе 0.50, но оба действия включаются с начального момента; это не пороговые реакции E7. E5 содержит 7,680 новых тестовых симуляций.

Серия со структурой E2/E4 включает последовательную и параллельную топологии, четыре мира и шесть независимых прогонов. В параллельной истинной модели после A следуют параллельные B/C с общим пулом из трёх исполнителей, AND-объединение, D и возможное повторение всего совместного блока. Признак сложности в этой серии отсутствует. Интенсивности обучения: 0.25/0.40/0.60/0.80; валидации: 0.35/0.70; поступления идут до момента 300 с последующим завершением всех задач. Inductive Miner в PM4Py 2.7.23.12 выявляет структуру; общие идентификаторы работников задают связанные пулы, число различных работников --- мощности. Композиционный исполнитель поддерживает последовательность, параллельные блоки и бинарные циклы. Одна известная функция возврата используется в выявленных циклах, что ограничивает интерпретацию переобобщения. Prosimos 2.1.0 исполняет выявленную BPMN с согласованными статическими параметрами. Базы последовательного/параллельного теста равны 0.50/0.40, всплеск 1.20 действует на $[200,400)$, горизонт 1,200, шесть повторов: 3,456 новых симуляций.

В прежних проверках ошибка траектории равна среднеквадратичному отклонению средних траекторий WIP модели и истинной системы, делённому на $\max(1,\text{средний WIP истины})$. Она вычисляется для каждого независимого обучающего набора, затем усредняется. Для парных разностей строятся процентильные интервалы по 10,000 бутстреп-перевыборкам шести независимых прогонов, описательные и без поправки. Эти оценки характеризуют восстановление траектории, а не фактическую пользу вмешательства.

Исследование реальных журналов использует исходные версии 1 BPI 2012 (13,087 задач; 262,200 событий) и BPI 2017 (31,509 задач; 1,202,267 событий)~\cite{bpi12,bpi17}. Интервалы начала/возобновления сопоставляются с наблюдаемыми остановками; это календарные интервалы, а не прямое измерение трудозатрат. Хронологические группы задач делятся 60/20/20; метки, пересекающие границы, удаляются. Операция, известный при старте ресурс, контекст задачи, предшествующая история и календарь образуют контекстный базовый метод. Нагрузка добавляет активные интервалы, открытые рабочие элементы и недавние старты глобально и по операциям/ресурсам. Неизвестный при старте ресурс остаётся неизвестным; восстановление состояния не зависит от будущего суффикса. В качестве сильного базового метода используется гистограммный градиентный бустинг. Цели: длительность $\log(1+\mathrm{seconds})$ и следующий наблюдаемый рабочий старт при условии продолжения, а не решение о выходе. Для парных потерь используются 2,000 бутстреп-перевыборок недельных блоков, шесть и шестнадцать тестовых недель. Проверяется прогнозная связь, а не политика в реальной системе.

\section{Результаты}
\subsection{Подготовка и включение реакций до кризиса}
Все 48 задач основной серии прошли проверку полноты, хешей завершения и исходных траекторий. Ни один из 528 калибровочных и валидационных эпизодов не достигает критерия тяжёлого события. Максимальный WIP при калибровке равен 16, 61, 19 и 111 для фиксированного мира, обслуживания, возвратов и совместного механизма соответственно. Последнее значение включает краткое превышение 100, поэтому <<до кризиса>> означает отсутствие заданного длительного события, а не постоянно низкую нагрузку. Ни один эпизод не исключён.

Таблица~\ref{tab:e7cost} показывает все механизмы и методы сравнения. В основном условии с возвратами планирование с учётом нагрузки выбирает сокращение возвратов при пороге 20 в 9 из 12 обучений, добавление мощности при пороге 20 --- в остальных 3. Средняя фактическая стоимость снижается с 92,832.6 у статического планирования до 57,351.8, то есть на 38.2\%. Сравнение с добавлением мощности при пороге 20 даёт снижение на 38.1\%. Оба скорректированных интервала в табл.~\ref{tab:e7contrast} исключают ноль: зафиксированный основной критерий выполнен. Фиксированное сокращение возвратов при пороге 20 даёт стоимость 45,313.5: результат модели численно хуже, а парный интервал включает ноль. Это не доказывает ни равенства, ни превосходства над данным правилом. Выбор компонент по прогнозной точности сохраняет зависимость возвратов от нагрузки лишь в 2 из 12 обучений и выбирает сокращение возвратов лишь в 1; его средняя стоимость равна 88,834.8.

\begin{table}[t]
\caption{Средняя фактическая стоимость E7 в тысячах условных единиц времени задачи, по 12 независимым обучающим наборам на механизм и всем тестовым условиям. <<Отбор>> означает выбор компонент по валидации.}\label{tab:e7cost}
\centering\small\setlength{\tabcolsep}{3pt}
\begin{tabular}{lrrrrrr}\toprule
Механизм & \shortstack{Без\\действия} & \shortstack{Статич.\\модель} & \shortstack{С\\нагрузкой} & \shortstack{С\\отбором} & \shortstack{Мощность\\при 20} & \shortstack{Возвраты\\при 20}\\\midrule
Фиксированный & 23.6 & 18.0 & 19.9 & 18.0 & 17.7 & 22.6 \\
Обслуживание & 211.7 & 111.0 & 111.0 & 111.0 & 111.0 & 202.7 \\
Возвраты & 147.8 & 92.8 & 57.4 & 88.8 & 92.6 & 45.3 \\
Совместный & 375.2 & 309.2 & 288.3 & 292.5 & 309.2 & 284.1 \\
\bottomrule

\end{tabular}
\end{table}
\begin{table}[t]
\caption{Возвраты: стоимость планирования с нагрузкой минус метод сравнения, в тысячах. Основные интервалы имеют уровень 97.5\%; сравнение с сокращением возвратов описательное, 95\%. Отрицательная разность означает выигрыш модели с нагрузкой.}\label{tab:e7contrast}
\centering\small
\begin{tabular}{lrrr}\toprule
Метод сравнения & Разность & Парный интервал & Уровень\\\midrule
Статич. планирование & $-35.48$ & $[-51.56, -19.40]$ & 97.5\% \\
Мощность при 20 & $-35.27$ & $[-51.23, -19.32]$ & 97.5\% \\
Возвраты при 20 & $12.04$ & $[-1.87, 25.94]$ & 95\% \\
\bottomrule

\end{tabular}
\end{table}

Реакция включается раньше кризиса в парной траектории без действия во всех 289 опасных траекториях мира с возвратами; средний запас времени равен 108.0 единицы. Она предотвращает заданное событие в 93 траекториях; в 196 оно сохраняется. В 192 траекториях без всплеска нет ни срабатываний, ни кризисов. Рис.~\ref{fig:e7} показывает ускоренное восстановление и разброс, вызванный различиями выбранных правил. Это условные описательные измерения предупреждения, а не оценка калиброванного классификатора кризисов.

\begin{figure}[t]
\centering\includegraphics[width=\textwidth]{generated/e7_response_ru.pdf}
\caption{Реакции в мире с возвратами для двух отложенных сценариев со всплеском. Восемь траекторий усредняются внутри каждого из 12 обучений, затем усредняются обучения. Полоса модели с нагрузкой --- одно стандартное отклонение средних по обучениям, а не доверительный интервал. Серые области отмечают всплески; горизонталь --- WIP 100. Фиксированное правило сокращения возвратов показывает ограничение выбора действия по модели.}\label{fig:e7}
\end{figure}

Польза зависит от механизма. При влиянии нагрузки на обслуживание все планировщики выбирают добавление мощности при пороге 20 и дают одинаковую стоимость. При совместном механизме модель с нагрузкой снижает стоимость на 6.7\% относительно статического планирования и правила добавления мощности, однако частота тяжёлых событий растёт с 69.3\% до 71.0\%: минимизация стоимости очереди не равна предотвращению заданного события. Без обратной связи модель с нагрузкой увеличивает стоимость на 10.9\% относительно статической, с описательной парной разностью 1,959.3 и интервалом $[219.5,3699.1]$. По четырём механизмам частота срабатываний без всплеска равна 0.0\%, 14.1\%, 0.0\% и 28.6\%; соответствующая частота кризисов без действий --- 0.0\%, 2.1\%, 0.0\% и 19.3\%. Неустойчивость базовой работы объясняет часть, но не все срабатывания без всплеска.

Анализ с равными весами механизмов при исходной стоимости даёт выигрыш модели с нагрузкой: нормированные разности равны $-0.0536$ относительно статического планирования, $-0.0503$ относительно добавления мощности при пороге 20 и $-0.1137$ относительно сокращения возвратов при пороге 20. Это разности отношений стоимости внутри миров, а не проценты снижения суммарной ненормированной стоимости. Таблица~\ref{tab:e7sensitivity} показывает описательные интервалы и все заранее заданные множители стоимости. При удвоенной стоимости действий интервалы относительно статического планирования и добавления мощности включают ноль; выигрыш относительно сокращения возвратов остаётся пограничным. Таким образом, приспособление реакции к механизму подтверждено для заданной задачи принятия решения, а не независимо от предположений о стоимости.

\begin{table}[t]
\caption{Нормированные разности стоимости по механизмам: планирование с нагрузкой минус метод сравнения, с описательными 95\%-ными интервалами. Для каждого множителя правила заново выбираются по стоимости планирования.}\label{tab:e7sensitivity}
\centering\small\setlength{\tabcolsep}{3pt}
\begin{tabular}{lrrr}\toprule
Множитель & Статич. планирование & Мощность при 20 & Возвраты при 20\\\midrule
0.5 & \shortstack{$-0.0467$\\$[-0.0772,-0.0161]$} & \shortstack{$-0.0467$\\$[-0.0772,-0.0161]$} & \shortstack{$-0.1201$\\$[-0.1501,-0.0901]$} \\
1.0 & \shortstack{$-0.0536$\\$[-0.0824,-0.0247]$} & \shortstack{$-0.0503$\\$[-0.0790,-0.0217]$} & \shortstack{$-0.1137$\\$[-0.1420,-0.0855]$} \\
2.0 & \shortstack{$-0.0136$\\$[-0.0618,0.0346]$} & \shortstack{$-0.0037$\\$[-0.0510,0.0436]$} & \shortstack{$-0.0471$\\$[-0.0939,-0.0002]$} \\
\bottomrule

\end{tabular}
\end{table}


\subsection{Вспомогательные свидетельства и ошибки переноса}
Модели предыдущей серии E1, обученные на широком диапазоне нагрузки, правильно выбирают пару активных механизмов во всех 24 наборах. Сравнение E5 с независимой случайностью улучшает эмпирический статический базовый метод в трёх механизмах при базе 0.65 со всплеском (табл.~\ref{tab:e5}). Снижения ошибки средней траектории составляют 76.2\%, 88.1\% и 64.4\%; все три парных описательных интервала лежат ниже нуля. Для фиксированного мира интервал включает ноль. Эти числа относятся к ошибке траектории в указанном условии, а не к стоимости E7 или универсальному ранжированию.
\begin{table}[t]
\caption{E5: всплеск при базе 0.65. Средняя NRMSE траектории WIP по шести независимым обучающим наборам. Разность: выбранная модель минус эмпирическая статическая; описательный 95\%-ный интервал.}\label{tab:e5}
\centering\small\setlength{\tabcolsep}{4pt}
\begin{tabular}{lrrrrl}\toprule
Механизм & Эмпирич. & Выбранная & Нейросеть & Разность & Интервал\\\midrule
Без зависимости & 0.0942 & 0.1067 & 0.2044 & $+0.0126$ & $[-0.0084,+0.0381]$ \\
Обслуживание & 0.1745 & 0.0416 & 0.0499 & $-0.1329$ & $[-0.1684,-0.1000]$ \\
Возвраты & 0.7733 & 0.0922 & 0.0750 & $-0.6811$ & $[-0.7488,-0.6118]$ \\
Оба механизма & 0.0920 & 0.0328 & 0.0475 & $-0.0592$ & $[-0.0693,-0.0496]$ \\
\bottomrule

\end{tabular}
\end{table}

При базе 0.50 истинный мир с возвратами заканчивает с 1.46 незавершёнными задачами при постоянном потоке и с 109.77 после всплеска, хотя интенсивность вернулась к базе за 1,500 единиц времени до конца. Обратная связь обслуживания также оставляет длительную очередь. Совместный механизм может быть неустойчив и без всплеска. Наблюдение подтверждает эффект затора на конечном горизонте, но не стационарный фазовый переход.

При выявленной структуре выбранное поведение улучшает согласованный Prosimos для последовательного обслуживания, последовательного совместного механизма и параллельного обслуживания (табл.~\ref{tab:structure}). Парные разности NRMSE равны $-0.1987$, $-0.7166$ и $-0.0998$, описательные интервалы исключают ноль. Параллельные возвраты и совместный механизм, напротив, ухудшаются на 0.2189 и 0.2019. Выявление допускает независимые повторения там, где истинный процесс повторяет совместный блок. Token fitness равен 1.0 для обеих топологий, но средний precision параллельной модели равен 0.4945. Таким образом, высокая полнота воспроизведения не гарантирует сохранения нагрузки выявленными циклами. Последовательные возвраты также выбираются статическими во всех прогонах при более узком диапазоне наблюдений.
\begin{table}[t]
\caption{Все условия со всплеском и выявленной структурой; средняя NRMSE по шести независимым обучающим наборам. D: выявленная структура; G: заданная.}\label{tab:structure}
\centering\small\setlength{\tabcolsep}{4pt}
\begin{tabular}{lrrrr}\toprule
Топология / механизм & D статич. & D выбранная & Prosimos & G выбранная\\\midrule
Послед. / Без зависимости & 0.1983 & 0.1874 & 0.2133 & 0.1850 \\
Послед. / Обслуживание & 0.9207 & 0.7219 & 0.9206 & 0.7119 \\
Послед. / Возвраты & 1.2318 & 1.2295 & 1.2351 & 1.2318 \\
Послед. / Оба механизма & 0.9779 & 0.2635 & 0.9801 & 0.2743 \\
Паралл. / Без зависимости & 0.0768 & 0.0914 & 0.0609 & 0.0800 \\
Паралл. / Обслуживание & 0.1968 & 0.0904 & 0.1901 & 0.0703 \\
Паралл. / Возвраты & 0.6397 & 0.8500 & 0.6312 & 0.4988 \\
Паралл. / Оба механизма & 0.1071 & 0.3141 & 0.1123 & 0.1559 \\
\bottomrule

\end{tabular}
\end{table}

На BPI 2017 MAE длительности у бустинга уменьшается с 1.188915 до 1.179020 при добавлении нагрузки: относительное улучшение 0.83\%, парный интервал $[-0.014642,-0.006249]$. Таблица~\ref{tab:real} сохраняет все пары журнала и цели. Длительность BPI 2012 почти не меняется, ошибка его следующего старта ухудшается, а интервал следующего старта BPI 2017 включает ноль. Нагрузка полезна для определённой цели; наблюдательные проверки не подтверждают реальное вмешательство.
\begin{table}[t]
\caption{Все реальные тесты бустинга. Длительность: MAE log(1+seconds); следующий старт: log loss. Интервалы описывают разности <<с нагрузкой минус контекст>>.}\label{tab:real}
\centering\small\setlength{\tabcolsep}{4pt}
\begin{tabular}{lrrrr}\toprule
Журнал / цель & Тест. строк & Контекст & С нагрузкой & Интервал разности\\\midrule
BPI12 / Длительность & 12,653 & 1.0308 & 1.0308 & $[-0.0051,+0.0033]$ \\
BPI12 / След. старт & 10,781 & 0.3585 & 0.3792 & $[+0.0022,+0.0436]$ \\
BPI17 / Длительность & 47,078 & 1.1889 & 1.1790 & $[-0.0146,-0.0062]$ \\
BPI17 / След. старт & 26,101 & 0.0815 & 0.0838 & $[-0.0005,+0.0058]$ \\
\bottomrule

\end{tabular}
\end{table}

Наконец, прежняя наблюдательная точность не проходит проверку ранжирования вмешательств. В совместном мире истинные конечные WIP для всплеска/B+1/половины возвратов равны 628.50/473.52/414.25, а выбранная модель прогнозирует 654.17/230.67/425.38 и предпочитает неверное действие. Диагностика координат после получения результата, E6, сохраняет знаменатель признака длительности B на обучающей мощности при увеличенной физической мощности. Прогноз становится 413.81, ошибка снижается примерно с 0.568 до 0.180, но ранжирование остаётся неверным. Диагностика добавляет 384 симуляции и не является независимо подтверждённым исправлением. E7 заранее фиксирует координаты признаков и проверяет фактическую выбранную реакцию, не считая наблюдательную точность достаточной.

\section{Обсуждение и ограничения}
Вклад в инженерию информационных систем --- процедура подготовки с явными свидетельствами на каждом этапе: восстановление рабочего состояния, моделирование поведения, гипотетическое нагрузочное испытание, выбор реакции и независимая оценка. Она связывает выявление процессов с проектированием и эксплуатацией систем. Карточка реакции может задавать наблюдаемый порог, требования к ресурсам, предполагаемый эффект действия, оценённую стоимость и проверенный диапазон нагрузок. Инструмент даёт вычислительную часть такой карточки; организационная осуществимость и внедрение остаются вне оценки.

Границы вывода существенны. E7 получает заданные последовательную структуру и смысл вмешательств. Вспомогательные эксперименты проверяют выявление структуры, но не повторяют сквозную политику E7 на выявленных параллельных процессах. Реальные журналы проверяют связь при неполных сведениях о жизненном цикле и ресурсах. Ни синтетический успех, ни небольшой реальный выигрыш в точности прогноза не идентифицируют причинный эффект вмешательства в организации. При внедрении нужно явно учитывать человеческое решение, наличие ресурсов, задержку действия и нормативные ограничения.

Ещё одна граница --- диапазон нагрузки. Спокойная валидация может поддерживать краткосрочный прогноз, но не выявлять зависимость, важную при будущем всплеске. Наоборот, неограниченная модель нагрузки способна экстраполировать шум в неправдоподобный кризис. Аудит калибровки и сравнения со статическим и выбранным вариантами показывают этот компромисс, но не решают общую задачу идентификации вне наблюдаемого распределения. Если порядок действий зависит от ненаблюдавшихся состояний, могут потребоваться дополнительные контролируемые измерения или безопасные нагрузочные пробы.

Двенадцать независимых обучающих наборов всё ещё дают грубую оценку неопределённости, а четыре синтетических мира не являются выборкой организаций. Основные интервалы корректируют два заранее заданных сравнения; остальные описательные. Равные веса сценариев и условная стоимость действий задают задачу принятия решения, а не наблюдаемые частоты инцидентов или цены. Мы сохраняем чувствительность к стоимости и равенство простым политикам, чтобы не завышать ценность обучения. Наблюдения с шагом пять, фиксированная длительность кризиса 200, известная эффективность действия и конечный горизонт ограничивают интерпретацию предотвращения и запаса времени. Постоянная устойчивость, надёжность для редких событий и перенос на пропуски событий либо смену состава задач не установлены.

Предыдущие серии содержат 11,952 новые тестовые симуляции отдельно от E7 и исключённого пилота. Число симуляций характеризует объём вычислений; неопределённость использует независимые обучения или хронологические блоки. Все неблагоприятные структурные и реальные сравнения сохраняются. Мы не заявляем превосходства над незапущенными полными системами выявления симуляций или причинного предписывающего мониторинга.

\section{Дополнительный диагностический протокол}
Следующие проверки были заданы после получения основных результатов. Исходные правила остаются зафиксированными; сохранены все 48 выборок и все заданные диагностические условия. Интервалы описательные, парные 95\%, по независимым обучающим выборкам.

\begin{table}[p]\caption{Средняя стоимость в тысячах для каждого механизма и диагностического сравнения.}\centering\small\setlength{\tabcolsep}{3pt}\begin{tabular}{p{.40\textwidth}rrrr}\toprule
Метод & Пост. & Длительность & Возвраты & Совместное\\\midrule
Без действия & 23.6 & 211.7 & 147.8 & 375.2\\
Статическое планирование & 18.0 & 111.0 & 92.8 & 309.2\\
Отбор по прогнозу & 18.0 & 111.0 & 88.8 & 292.5\\
Правило по силе зависимости & 18.9 & 111.0 & 73.0 & 309.2\\
Сохранение нагрузочных компонент & 19.9 & 111.0 & 57.4 & 288.3\\
Ресурс при 20 & 17.7 & 111.0 & 92.6 & 309.2\\
Возвраты при 20 & 22.6 & 202.7 & 45.3 & 284.1\\
Известный генератор & 17.8 & 111.0 & 45.3 & 284.1\\
Ретроспективная граница & 17.7 & 111.0 & 45.3 & 284.1\\
\bottomrule\end{tabular}\end{table}
\begin{table}[p]\caption{Превышение стоимости относительно планировщика с известным генератором: среднее и 95\% интервал, в тысячах. Отрицательные значения возможны из-за конечной выборки и иных сценариев планирования.}\centering\small\setlength{\tabcolsep}{3pt}\begin{tabular}{llr}\toprule
Механизм & Метод & Превышение [интервал]\\\midrule
Постоянное & Статическое планирование & 0.17 [-0.08, 0.42]\\
Постоянное & Отбор по прогнозу & 0.17 [-0.08, 0.42]\\
Постоянное & Правило по силе зависимости & 1.09 [-0.19, 2.36]\\
Постоянное & Сохранение нагрузочных компонент & 2.13 [0.41, 3.85]\\
Длительность & Статическое планирование & 0.00 [0.00, 0.00]\\
Длительность & Отбор по прогнозу & 0.00 [0.00, 0.00]\\
Длительность & Правило по силе зависимости & 0.00 [0.00, 0.00]\\
Длительность & Сохранение нагрузочных компонент & 0.00 [0.00, 0.00]\\
Возвраты & Статическое планирование & 47.52 [45.70, 49.34]\\
Возвраты & Отбор по прогнозу & 43.52 [34.63, 52.42]\\
Возвраты & Правило по силе зависимости & 27.66 [12.12, 43.20]\\
Возвраты & Сохранение нагрузочных компонент & 12.04 [-1.87, 25.94]\\
Совместное & Статическое планирование & 25.05 [23.43, 26.68]\\
Совместное & Отбор по прогнозу & 8.32 [0.44, 16.19]\\
Совместное & Правило по силе зависимости & 25.05 [23.43, 26.68]\\
Совместное & Сохранение нагрузочных компонент & 4.20 [-2.05, 10.45]\\
\bottomrule\end{tabular}\end{table}
\begin{table}[p]\caption{Эффект вмешательства: решения зафиксированы при предполагаемом сокращении 50\%; показаны фактические сокращения. Средние стоимости в тысячах.}\centering\small\setlength{\tabcolsep}{3pt}\begin{tabular}{llrrrrr}\toprule
Механизм & Эффект & Стат. & Отбор & Нагрузка & Ресурс & Возвраты\\\midrule
Постоянное & 30\% & 18.0 & 18.0 & 20.1 & 17.7 & 23.2\\
Постоянное & 50\% & 18.0 & 18.0 & 19.9 & 17.7 & 22.6\\
Постоянное & 70\% & 18.0 & 18.0 & 19.7 & 17.7 & 22.0\\
Длительность & 30\% & 111.0 & 111.0 & 111.0 & 111.0 & 206.5\\
Длительность & 50\% & 111.0 & 111.0 & 111.0 & 111.0 & 202.7\\
Длительность & 70\% & 111.0 & 111.0 & 111.0 & 111.0 & 199.1\\
Возвраты & 30\% & 92.8 & 91.1 & 77.1 & 92.6 & 71.7\\
Возвраты & 50\% & 92.8 & 88.8 & 57.4 & 92.6 & 45.3\\
Возвраты & 70\% & 92.8 & 87.8 & 47.6 & 92.6 & 32.2\\
Совместное & 30\% & 309.2 & 316.8 & 319.2 & 309.2 & 321.1\\
Совместное & 50\% & 309.2 & 292.5 & 288.3 & 309.2 & 284.1\\
Совместное & 70\% & 309.2 & 268.1 & 258.0 & 309.2 & 248.1\\
\bottomrule\end{tabular}\end{table}
\begin{table}[p]\caption{Чувствительность к смеси: средняя нормированная разность нагрузочного метода и сравнения; отрицательное значение благоприятно. Меняется вес постоянного механизма, остальные три делят оставшийся вес поровну.}\centering\small\setlength{\tabcolsep}{3pt}\begin{tabular}{p{.37\textwidth}rrrr}\toprule
Сравнение & 0 & 0.25 & 0.5 & 1\\\midrule
Статическое планирование & -0.099 & -0.054 & -0.008 & 0.083\\
Отбор по прогнозу & -0.075 & -0.036 & 0.004 & 0.083\\
Ресурс при 20 & -0.099 & -0.050 & -0.002 & 0.095\\
Возвраты при 20 & -0.114 & -0.114 & -0.113 & -0.113\\
Правило по силе зависимости & -0.054 & -0.030 & -0.005 & 0.044\\
Известный генератор & 0.030 & 0.045 & 0.060 & 0.090\\
\bottomrule\end{tabular}\end{table}
\begin{table}[p]\caption{Чувствительность критерия события для возвратов: процент траекторий, достигших критерия; исходные зафиксированные правила.}\centering\small\setlength{\tabcolsep}{3pt}\begin{tabular}{rrrrrr}\toprule
WIP & Длительность & Без действия & Стат. & Нагрузка & Возвраты\\\midrule
80 & 100 & 61.6 & 50.5 & 39.9 & 37.5\\
80 & 200 & 59.0 & 43.6 & 37.3 & 34.2\\
80 & 300 & 55.6 & 38.2 & 34.5 & 33.3\\
100 & 100 & 54.2 & 39.2 & 35.6 & 33.7\\
100 & 200 & 50.2 & 35.8 & 34.0 & 33.3\\
100 & 300 & 46.9 & 34.4 & 33.3 & 32.5\\
120 & 100 & 43.9 & 34.4 & 34.0 & 33.3\\
120 & 200 & 42.7 & 34.0 & 33.2 & 32.1\\
120 & 300 & 41.3 & 33.9 & 30.4 & 29.0\\
\bottomrule\end{tabular}\end{table}
\begin{table}[p]\caption{Чувствительность стоимости для возвратов: горизонт и вес конечного остатка. Средние стоимости в тысячах; реакции повторно не выбираются.}\centering\small\setlength{\tabcolsep}{3pt}\begin{tabular}{rrrrrr}\toprule
Горизонт & Вес & Стат. & Отбор & Нагрузка & Возвраты\\\midrule
1400 & 0 & 77.0 & 74.2 & 52.2 & 43.7\\
1400 & 100 & 81.9 & 78.7 & 53.8 & 44.3\\
1400 & 200 & 86.8 & 83.2 & 55.4 & 44.9\\
1600 & 0 & 85.9 & 82.4 & 55.1 & 44.6\\
1600 & 100 & 89.7 & 85.9 & 56.3 & 44.9\\
1600 & 200 & 93.5 & 89.5 & 57.4 & 45.3\\
1800 & 0 & 92.8 & 88.8 & 57.4 & 45.3\\
1800 & 100 & 95.8 & 91.6 & 58.4 & 45.6\\
1800 & 200 & 98.7 & 94.3 & 59.4 & 45.9\\
\bottomrule\end{tabular}\end{table}
Архив содержит полный выбор действий по каждой выборке, распределения порогов, исходные условия эффективности и чувствительности, выявленные сети, проверки компилятора и машиночитаемую карточку реакции. 4032 новых запусков планирования с известным генератором; 9156 новых запусков с изменённым эффектом и 4668 точно повторно использованных траекторий без активации. Повторное использование допустимо, поскольку эффект меняет только события после активации, которая на этих траекториях не происходит.
\clearpage
\begin{table}[p]\caption{Числа включений и предотвращений при зафиксированных нагрузочных реакциях. Последний столбец считает включения на траекториях без тяжёлого события в сопоставленном запуске без действия; это отличается от включения в спокойном контроле.}\centering\small\setlength{\tabcolsep}{3pt}\begin{tabular}{lrrrr}\toprule
Механизм & Риск & Ранние & Предотвращено & \shortstack{Включение без\\события в базе}\\\midrule
Постоянное & 117 & 109 & 74 & 239\\
Длительность & 382 & 382 & 138 & 29\\
Возвраты & 289 & 289 & 93 & 95\\
Совместное & 421 & 421 & 12 & 18\\
\bottomrule\end{tabular}\end{table}\clearpage

\section{Последовательный процесс: чувствительность селектора}
Диагностика P0 повторно использует все сохранённые тестовые действия, изменяя порог отбора возвратов по сетке $0,0.001,0.0025,0.005,0.01$. Компонент сохраняется соответственно в $8,5,2,2,0$ из 12 routing-обучений; средние стоимости составляют $80.8,84.9,88.8,88.8,92.8$ тыс. Стоимость retained-load остаётся $57.4$ тыс. Следовательно, наблюдаемое различие не ограничивается исходным порогом $0.005$. При порогах отбора длительности $0,0.0025,0.005,0.01,0.02$ стоимость селектора равна $85.1$ тыс. при нулевом пороге и $88.8$ при остальных. Это диагностика на прежнем независимом тестовом каталоге, а не новые подтверждающие наборы. Сохранение только зависимости длительности стоит $69.4$ тыс., только возвратов — $80.8$ тыс.; оба компонента дают минимальную среднюю стоимость среди этих абляций.

\section{Исполнимая маршрутизация и аудит разработки}
Структурное исследование использует новый зафиксированный протокол и новые случайные потоки обучения, валидации, планирования и тестирования. Прежняя независимая реализация сохранена в версионированном исследовательском архиве. Её шесть отдельных сравнений при routing feedback и известной структуре не установили положительного эффекта после поправки на множественные проверки; описательная нормированная разность с равными весами процессов составила $-0.090$ при заданной структуре и $0.001$ при обнаруженной. Эти исходы не объединяются с новыми основными сравнениями.

До новой фиксации аналитический аудит выявил две ошибки исполнителя. Два дерева с наблюдаемым языком $(BD)^+$ размещали решение о возврате в разных точках. При фиксированной вероятности возврата $0.2$ прежний исполнитель давал 1.250 и 2.027 прохода на экземпляр; уменьшение вероятности вдвое давало 1.1085 и неизменные 2.027. Проверка содержит 2000 независимо сгенерированных экземпляров. Преобразование общего продолжения переносит выбор в наблюдаемую точку решения и обеспечивает точное совпадение двух массивов траекторий с вмешательством и без него. Ожидаемые числа проходов равны $1/0.8$ и $1/0.9$. Отдельная проверка XOR с условием выявила 999 нарушений среди 2000 экземпляров: имитация обученной модели игнорировала явные условия по классу. Исправленный исполнитель не допускает нарушений. Эти проверки устанавливают свойства исполнения, а не величину эффекта по стоимости политики.

Отдельный аудит из 92 запусков охватывает каждый процесс, механизм обратной связи и тип вмешательства. Отсчёты истинного генератора, записи жизненного цикла, завершения и времена активации идентичны до и после исправления. Уравнения генерации, мощности, стоимости, условия поступлений и масштабы нагрузки не меняются. Технические пилоты проверяют все шесть обнаруженных структур и эквивалентность планировочных оценок при известной и обнаруженной структуре; они исключены из статистических выводов.

Процедура обнаружения сравнивает стандартный Inductive Miner с кандидатом, который отдаёт приоритет существующему правилу строгого цикла и выводит детерминированные условия по наблюдаемым атрибутам. Допускаются только кандидаты, которым соответствуют все обучающие трассы по оптимальным выравниваниям и которые проходят проверки компилятора; затем максимизируется обучающая precision по token replay. Глобальная alignment fitness может быть чуть ниже единицы из-за ненулевой стоимости невидимых переходов; допуск использует долю соответствующих трасс. Кандидат не получает исходное дерево. Эквивалентность наблюдаемых языков проверяется только затем как диагностика на синтетических данных. Сохранены все кандидатные модели, условия, результаты проверки соответствия и квитанции подготовки.

Исполненная реализация ограничивает модель 30 наблюдаемыми метками, 62 позициями, 120 переходами и 32 ресурсными пулами при бинарных наблюдаемых атрибутах класса и процесса. Неподдерживаемые развилки и небезопасные маркировки отклоняются. Это ограничения реализации, а не эмпирическое утверждение о масштабируемости.

Дополнительная техническая проверка использует только исключённые из оценки пилоты. Общий объект конфигурации содержит поля истинного генератора, хотя подготовка на обученной модели их не использует. Подмена истинного дерева и генеративных коэффициентов, а также удаление топологии из входов обучения поведения оставляют обученные параметры, выбранную обнаруженную структуру и исполнение обученной модели полностью неизменными во всех шести пилотах. Ресурсные отображения, потенциальные точки возврата и наблюдаемые очереди остаются инженерными входами.

\clearpage

\section{Независимое структурное расширение}

Расширение содержит 288 независимых обучающих наборов и 3168 проверенных калибровочных/валидационных эпизодов. Предкризисный критерий нарушен в 1 эпизодах; процесс, механизм, обучение и номер эпизода записаны в \texttt{calibration\_audits.json}, результаты сохранены. Основные выводы при routing feedback требуют полного предкризисного покрытия именно в этом механизме. Известная и обнаруженная структуры используют одинаковое обученное поведение и заданные ресурсные отображения. Полные сравнения методов, действий и компонентов находятся в \texttt{analysis/full\_summary.csv}.

\subsection{Стоимость по всем структурам и механизмам}

Таблицы показывают среднее отношений стоимости каждого обученного варианта к соответствующей стоимости без вмешательства, равной 1. Regret — разность retained-load и планировщика с генератором в этой шкале. Частота тяжёлых событий относится к retained-load. K означает заданную структуру, D — выбранную обнаруженную, L — стандартный IM с тем же исправленным исполнителем. В каждой строке 12 независимых обучений; эвристика и генератор одинаковы в всех структурных вариантах.

\begin{table}[htbp]\caption{P1: нормированная стоимость и частота событий.}\centering\small\setlength{\tabcolsep}{3pt}

\begin{tabular}{llrrrrrrr}\toprule

Model & Mechanism & Static & Pred. & Ret. & Heur. & Gen. & Regret & Event\\\midrule

K & fixed & 0.816 & 0.816 & 0.810 & 0.805 & 0.810 & -0.001 & 0.076\\

K & service & 0.594 & 0.590 & 0.608 & 0.611 & 0.579 & 0.029 & 0.368\\

K & routing & 0.584 & 0.555 & 0.491 & 0.499 & 0.429 & 0.062 & 0.365\\

K & combined & 0.734 & 0.675 & 0.637 & 0.625 & 0.592 & 0.045 & 0.649\\

D & fixed & 0.816 & 0.816 & 0.810 & 0.805 & 0.810 & -0.001 & 0.076\\

D & service & 0.594 & 0.590 & 0.608 & 0.611 & 0.579 & 0.029 & 0.368\\

D & routing & 0.584 & 0.555 & 0.491 & 0.499 & 0.429 & 0.062 & 0.365\\

D & combined & 0.734 & 0.675 & 0.637 & 0.625 & 0.592 & 0.045 & 0.649\\

L & fixed & 0.781 & 0.781 & 0.767 & 0.805 & 0.810 & -0.043 & 0.036\\

L & service & 0.581 & 0.580 & 0.579 & 0.611 & 0.579 & 0.000 & 0.359\\

L & routing & 0.570 & 0.570 & 0.566 & 0.499 & 0.429 & 0.137 & 0.375\\

L & combined & 0.728 & 0.728 & 0.725 & 0.625 & 0.592 & 0.133 & 0.656\\

\bottomrule\end{tabular}\end{table}

\begin{table}[htbp]\caption{P2: нормированная стоимость и частота событий.}\centering\small\setlength{\tabcolsep}{3pt}

\begin{tabular}{llrrrrrrr}\toprule

Model & Mechanism & Static & Pred. & Ret. & Heur. & Gen. & Regret & Event\\\midrule

K & fixed & 0.958 & 0.955 & 0.966 & 1.096 & 0.965 & 0.002 & 0.082\\

K & service & 0.646 & 0.646 & 0.830 & 0.850 & 0.626 & 0.204 & 0.340\\

K & routing & 0.688 & 0.658 & 0.561 & 0.648 & 0.435 & 0.126 & 0.345\\

K & combined & 0.750 & 0.749 & 0.575 & 0.625 & 0.516 & 0.059 & 0.502\\

D & fixed & 0.958 & 0.955 & 0.966 & 1.096 & 0.965 & 0.002 & 0.082\\

D & service & 0.646 & 0.646 & 0.830 & 0.850 & 0.626 & 0.204 & 0.340\\

D & routing & 0.688 & 0.658 & 0.561 & 0.648 & 0.435 & 0.126 & 0.345\\

D & combined & 0.750 & 0.749 & 0.575 & 0.625 & 0.516 & 0.059 & 0.502\\

L & fixed & 0.923 & 0.923 & 1.087 & 1.096 & 0.965 & 0.122 & 0.102\\

L & service & 0.632 & 0.632 & 0.880 & 0.850 & 0.626 & 0.254 & 0.344\\

L & routing & 0.781 & 0.781 & 0.898 & 0.648 & 0.435 & 0.463 & 0.378\\

L & combined & 0.744 & 0.753 & 0.787 & 0.625 & 0.516 & 0.271 & 0.562\\

\bottomrule\end{tabular}\end{table}

\begin{table}[htbp]\caption{P3: нормированная стоимость и частота событий.}\centering\small\setlength{\tabcolsep}{3pt}

\begin{tabular}{llrrrrrrr}\toprule

Model & Mechanism & Static & Pred. & Ret. & Heur. & Gen. & Regret & Event\\\midrule

K & fixed & 0.922 & 0.922 & 0.925 & 0.908 & 0.911 & 0.014 & 0.293\\

K & service & 0.529 & 0.529 & 0.738 & 0.601 & 0.528 & 0.210 & 0.568\\

K & routing & 0.659 & 0.659 & 0.613 & 0.635 & 0.522 & 0.092 & 0.642\\

K & combined & 0.878 & 0.875 & 0.862 & 0.881 & 0.837 & 0.025 & 0.667\\

D & fixed & 0.922 & 0.922 & 0.925 & 0.908 & 0.911 & 0.014 & 0.293\\

D & service & 0.529 & 0.529 & 0.738 & 0.601 & 0.528 & 0.210 & 0.568\\

D & routing & 0.659 & 0.659 & 0.613 & 0.635 & 0.522 & 0.092 & 0.642\\

D & combined & 0.878 & 0.875 & 0.862 & 0.881 & 0.837 & 0.025 & 0.667\\

L & fixed & 0.922 & 0.922 & 0.925 & 0.908 & 0.911 & 0.014 & 0.293\\

L & service & 0.529 & 0.529 & 0.738 & 0.601 & 0.528 & 0.210 & 0.568\\

L & routing & 0.659 & 0.659 & 0.613 & 0.635 & 0.522 & 0.092 & 0.642\\

L & combined & 0.878 & 0.875 & 0.862 & 0.881 & 0.837 & 0.025 & 0.667\\

\bottomrule\end{tabular}\end{table}

\begin{table}[htbp]\caption{P4: нормированная стоимость и частота событий.}\centering\small\setlength{\tabcolsep}{3pt}

\begin{tabular}{llrrrrrrr}\toprule

Model & Mechanism & Static & Pred. & Ret. & Heur. & Gen. & Regret & Event\\\midrule

K & fixed & 0.722 & 0.743 & 0.807 & 0.803 & 0.722 & 0.086 & 0.240\\

K & service & 0.543 & 0.543 & 0.610 & 0.610 & 0.542 & 0.068 & 0.557\\

K & routing & 0.518 & 0.527 & 0.615 & 0.623 & 0.513 & 0.102 & 0.495\\

K & combined & 0.798 & 0.818 & 0.822 & 0.826 & 0.798 & 0.024 & 0.667\\

D & fixed & 0.722 & 0.743 & 0.807 & 0.803 & 0.722 & 0.086 & 0.240\\

D & service & 0.543 & 0.543 & 0.610 & 0.610 & 0.542 & 0.068 & 0.557\\

D & routing & 0.518 & 0.527 & 0.615 & 0.623 & 0.513 & 0.102 & 0.495\\

D & combined & 0.798 & 0.818 & 0.822 & 0.826 & 0.798 & 0.024 & 0.667\\

L & fixed & 0.722 & 0.743 & 0.807 & 0.803 & 0.722 & 0.086 & 0.240\\

L & service & 0.543 & 0.543 & 0.610 & 0.610 & 0.542 & 0.068 & 0.557\\

L & routing & 0.518 & 0.527 & 0.615 & 0.623 & 0.513 & 0.102 & 0.495\\

L & combined & 0.798 & 0.818 & 0.822 & 0.826 & 0.798 & 0.024 & 0.667\\

\bottomrule\end{tabular}\end{table}

\begin{table}[htbp]\caption{P5: нормированная стоимость и частота событий.}\centering\small\setlength{\tabcolsep}{3pt}

\begin{tabular}{llrrrrrrr}\toprule

Model & Mechanism & Static & Pred. & Ret. & Heur. & Gen. & Regret & Event\\\midrule

K & fixed & 0.798 & 0.798 & 0.925 & 1.032 & 0.799 & 0.126 & 0.309\\

K & service & 0.350 & 0.350 & 0.353 & 0.512 & 0.348 & 0.005 & 0.470\\

K & routing & 1.030 & 0.828 & 0.571 & 0.467 & 0.415 & 0.155 & 0.382\\

K & combined & 1.139 & 1.043 & 1.018 & 1.018 & 0.979 & 0.039 & 0.674\\

D & fixed & 0.798 & 0.798 & 0.925 & 1.032 & 0.799 & 0.126 & 0.309\\

D & service & 0.350 & 0.350 & 0.353 & 0.512 & 0.348 & 0.005 & 0.470\\

D & routing & 1.030 & 0.828 & 0.571 & 0.467 & 0.415 & 0.155 & 0.382\\

D & combined & 1.139 & 1.043 & 1.018 & 1.018 & 0.979 & 0.039 & 0.674\\

L & fixed & 0.798 & 0.798 & 0.925 & 1.032 & 0.799 & 0.126 & 0.309\\

L & service & 0.350 & 0.350 & 0.353 & 0.512 & 0.348 & 0.005 & 0.470\\

L & routing & 1.030 & 0.828 & 0.571 & 0.467 & 0.415 & 0.155 & 0.382\\

L & combined & 1.139 & 1.043 & 1.018 & 1.018 & 0.979 & 0.039 & 0.674\\

\bottomrule\end{tabular}\end{table}

\begin{table}[htbp]\caption{P6: нормированная стоимость и частота событий.}\centering\small\setlength{\tabcolsep}{3pt}

\begin{tabular}{llrrrrrrr}\toprule

Model & Mechanism & Static & Pred. & Ret. & Heur. & Gen. & Regret & Event\\\midrule

K & fixed & 0.754 & 0.754 & 0.854 & 1.034 & 0.749 & 0.105 & 0.253\\

K & service & 0.496 & 0.496 & 0.671 & 0.878 & 0.493 & 0.178 & 0.464\\

K & routing & 0.573 & 0.537 & 0.404 & 0.574 & 0.351 & 0.054 & 0.382\\

K & combined & 0.744 & 0.680 & 0.598 & 0.656 & 0.585 & 0.013 & 0.667\\

D & fixed & 0.754 & 0.754 & 0.854 & 1.034 & 0.749 & 0.105 & 0.253\\

D & service & 0.496 & 0.496 & 0.671 & 0.878 & 0.493 & 0.178 & 0.464\\

D & routing & 0.573 & 0.537 & 0.404 & 0.574 & 0.351 & 0.054 & 0.382\\

D & combined & 0.744 & 0.680 & 0.598 & 0.656 & 0.585 & 0.013 & 0.667\\

L & fixed & 0.739 & 0.739 & 0.739 & 1.034 & 0.749 & -0.010 & 0.165\\

L & service & 0.493 & 0.493 & 0.493 & 0.878 & 0.493 & 0.000 & 0.365\\

L & routing & 0.567 & 0.567 & 0.567 & 0.574 & 0.351 & 0.217 & 0.493\\

L & combined & 0.739 & 0.739 & 0.739 & 0.656 & 0.585 & 0.154 & 0.667\\

\bottomrule\end{tabular}\end{table}

\clearpage

\subsection{Основные сравнения и проверка обнаружения структуры}

Два основных сравнения при нагрузочных возвратах усредняют парные разности стоимости, нормированной на no action, с равными весами шести процессов. Двусторонние интервалы 97.5\% используют независимую дисперсию по обучениям и приближение Саттертуэйта; ошибка семейства двух сравнений равна 0.05. H1 сравнивает сохранение зависимостей с предиктивным отбором при выбранной обнаруженной структуре; H2 сравнивает сохранение зависимостей при выбранной структуре и стандартном IM. Отрицательные разности соответствуют преимуществу предложенного варианта.

\begin{center}\small\begin{tabular}{lrrr}\toprule

Contrast & Difference & Lower & Upper\\\midrule

H1 & -0.0848 & -0.1353 & -0.0343\\

H2 & -0.0958 & -0.1205 & -0.0711\\

\bottomrule\end{tabular}\end{center}

Следующие шесть сравнений отдельных процессов при нагрузочных возвратах вторичны. Интервалы Бонферрони охватывают это вторичное семейство. Значения — абсолютные разности стоимости retained минус predictive, делённые на 1000, при выбранной обнаруженной структуре.

\begin{center}\small\begin{tabular}{lrrr}\toprule

Process & Difference & Lower & Upper\\\midrule

P1 & -1.595 & -3.226 & 0.037\\

P2 & -2.452 & -8.808 & 3.903\\

P3 & -4.998 & -11.930 & 1.934\\

P4 & 5.524 & 2.304 & 8.744\\

P5 & -15.714 & -33.446 & 2.019\\

P6 & -10.153 & -18.481 & -1.824\\

\bottomrule\end{tabular}\end{center}

\begin{center}\small\begin{tabular}{lrrrrrr}\toprule

Process & Fitness & Precision & Equal & Added & Missing & Same policy\\\midrule

P1 & 1.000 & 0.999 & 48/48 & 0/48 & 0/48 & 48/48\\

P2 & 1.000 & 1.000 & 48/48 & 0/48 & 0/48 & 48/48\\

P3 & 1.000 & 0.996 & 48/48 & 0/48 & 0/48 & 48/48\\

P4 & 1.000 & 1.000 & 48/48 & 0/48 & 0/48 & 48/48\\

P5 & 1.000 & 1.000 & 48/48 & 0/48 & 0/48 & 48/48\\

P6 & 1.000 & 1.000 & 48/48 & 0/48 & 0/48 & 48/48\\

\bottomrule\end{tabular}\end{center}

Проверка языков сравнивает видимые последовательности исполнения с фиксированным ограничением числа состояний; незавершённые проверки сохранены в \texttt{discovery\_summary.json}. Same policy считает полное совпадение действия и порога retained-load между известной и обнаруженной структурой во всех четырёх механизмах. Добавленные или пропущенные пути регистрируются как нарушения соответствия исходной структуре в \texttt{structural\_conformance.json}; их стоимость без исправления сохраняется для анализа ошибок. Этот признак оценки по известному оригиналу отделён от вычислительного отказа и не является реализованным правилом допуска политики.

\subsection{Чувствительность селектора}

Перебор использует сохранённые планировочные решения и независимые тестовые действия. Это диагностика; порог не выбирается по тестовой стоимости. Таблицы относятся к нагрузочным возвратам и известной структуре. D и R — число сохранённых компонентов из 12. Agree — совпадение типа действия с планировщиком по генератору; regret — абсолютная разность стоимости, делённая на 1000. Остальные механизмы и обнаруженные структуры включены в \texttt{selector\_summary.csv}.

\begingroup\small\setlength{\tabcolsep}{3pt}\begin{longtable}{llrrrrrr}\caption{Перебор селектора: return.}\\\toprule

Process & Threshold & D & R & Agree & Cost ratio & Regret & Event\\\midrule\endfirsthead

\toprule Process & Threshold & D & R & Agree & Cost ratio & Regret & Event\\\midrule\endhead

P1 & 0 & 0 & 4 & 0.250 & 0.542 & 2.829 & 0.377\\

P1 & 0.001 & 0 & 2 & 0.167 & 0.555 & 3.143 & 0.380\\

P1 & 0.0025 & 0 & 2 & 0.167 & 0.555 & 3.143 & 0.380\\

P1 & 0.005 & 0 & 2 & 0.167 & 0.555 & 3.143 & 0.380\\

P1 & 0.01 & 0 & 1 & 0.083 & 0.570 & 3.523 & 0.387\\

P2 & 0 & 0 & 9 & 0.583 & 0.617 & 4.574 & 0.359\\

P2 & 0.001 & 0 & 7 & 0.500 & 0.649 & 5.374 & 0.368\\

P2 & 0.0025 & 0 & 6 & 0.417 & 0.677 & 6.100 & 0.372\\

P2 & 0.005 & 0 & 3 & 0.417 & 0.658 & 5.655 & 0.370\\

P2 & 0.01 & 0 & 1 & 0.417 & 0.660 & 5.703 & 0.372\\

P3 & 0 & 0 & 6 & 0.250 & 0.646 & 13.541 & 0.642\\

P3 & 0.001 & 0 & 2 & 0.167 & 0.642 & 13.178 & 0.646\\

P3 & 0.0025 & 0 & 0 & 0.083 & 0.659 & 14.965 & 0.649\\

P3 & 0.005 & 0 & 0 & 0.083 & 0.659 & 14.965 & 0.649\\

P3 & 0.01 & 0 & 0 & 0.083 & 0.659 & 14.965 & 0.649\\

P4 & 0 & 0 & 7 & 0.500 & 0.576 & 3.909 & 0.462\\

P4 & 0.001 & 0 & 7 & 0.500 & 0.576 & 3.909 & 0.462\\

P4 & 0.0025 & 0 & 6 & 0.583 & 0.566 & 3.306 & 0.446\\

P4 & 0.005 & 0 & 1 & 0.917 & 0.527 & 0.856 & 0.422\\

P4 & 0.01 & 0 & 1 & 0.917 & 0.527 & 0.856 & 0.422\\

P5 & 0 & 0 & 8 & 0.500 & 0.726 & 19.055 & 0.429\\

P5 & 0.001 & 0 & 7 & 0.417 & 0.777 & 22.173 & 0.444\\

P5 & 0.0025 & 0 & 7 & 0.417 & 0.777 & 22.173 & 0.444\\

P5 & 0.005 & 0 & 6 & 0.333 & 0.828 & 25.348 & 0.458\\

P5 & 0.01 & 0 & 2 & 0.167 & 0.928 & 31.265 & 0.486\\

P6 & 0 & 0 & 5 & 0.333 & 0.499 & 11.188 & 0.453\\

P6 & 0.001 & 0 & 4 & 0.333 & 0.499 & 11.188 & 0.453\\

P6 & 0.0025 & 0 & 4 & 0.333 & 0.499 & 11.188 & 0.453\\

P6 & 0.005 & 0 & 2 & 0.167 & 0.537 & 14.175 & 0.483\\

P6 & 0.01 & 0 & 0 & 0.000 & 0.573 & 16.918 & 0.509\\

\bottomrule\end{longtable}\endgroup

\begingroup\small\setlength{\tabcolsep}{3pt}\begin{longtable}{llrrrrrr}\caption{Перебор селектора: duration.}\\\toprule

Process & Threshold & D & R & Agree & Cost ratio & Regret & Event\\\midrule\endfirsthead

\toprule Process & Threshold & D & R & Agree & Cost ratio & Regret & Event\\\midrule\endhead

P1 & 0 & 2 & 2 & 0.167 & 0.555 & 3.143 & 0.380\\

P1 & 0.0025 & 1 & 2 & 0.167 & 0.555 & 3.143 & 0.380\\

P1 & 0.005 & 0 & 2 & 0.167 & 0.555 & 3.143 & 0.380\\

P1 & 0.01 & 0 & 2 & 0.167 & 0.555 & 3.143 & 0.380\\

P1 & 0.02 & 0 & 2 & 0.167 & 0.555 & 3.143 & 0.380\\

P2 & 0 & 3 & 3 & 0.333 & 0.697 & 6.752 & 0.372\\

P2 & 0.0025 & 0 & 3 & 0.417 & 0.658 & 5.655 & 0.370\\

P2 & 0.005 & 0 & 3 & 0.417 & 0.658 & 5.655 & 0.370\\

P2 & 0.01 & 0 & 3 & 0.417 & 0.658 & 5.655 & 0.370\\

P2 & 0.02 & 0 & 3 & 0.417 & 0.658 & 5.655 & 0.370\\

P3 & 0 & 2 & 0 & 0.083 & 0.657 & 14.768 & 0.648\\

P3 & 0.0025 & 0 & 0 & 0.083 & 0.659 & 14.965 & 0.649\\

P3 & 0.005 & 0 & 0 & 0.083 & 0.659 & 14.965 & 0.649\\

P3 & 0.01 & 0 & 0 & 0.083 & 0.659 & 14.965 & 0.649\\

P3 & 0.02 & 0 & 0 & 0.083 & 0.659 & 14.965 & 0.649\\

P4 & 0 & 1 & 1 & 0.917 & 0.527 & 0.856 & 0.422\\

P4 & 0.0025 & 0 & 1 & 0.917 & 0.527 & 0.856 & 0.422\\

P4 & 0.005 & 0 & 1 & 0.917 & 0.527 & 0.856 & 0.422\\

P4 & 0.01 & 0 & 1 & 0.917 & 0.527 & 0.856 & 0.422\\

P4 & 0.02 & 0 & 1 & 0.917 & 0.527 & 0.856 & 0.422\\

P5 & 0 & 3 & 6 & 0.333 & 0.826 & 25.218 & 0.457\\

P5 & 0.0025 & 0 & 6 & 0.333 & 0.828 & 25.348 & 0.458\\

P5 & 0.005 & 0 & 6 & 0.333 & 0.828 & 25.348 & 0.458\\

P5 & 0.01 & 0 & 6 & 0.333 & 0.828 & 25.348 & 0.458\\

P5 & 0.02 & 0 & 6 & 0.333 & 0.828 & 25.348 & 0.458\\

P6 & 0 & 1 & 2 & 0.167 & 0.537 & 14.175 & 0.483\\

P6 & 0.0025 & 0 & 2 & 0.167 & 0.537 & 14.175 & 0.483\\

P6 & 0.005 & 0 & 2 & 0.167 & 0.537 & 14.175 & 0.483\\

P6 & 0.01 & 0 & 2 & 0.167 & 0.537 & 14.175 & 0.483\\

P6 & 0.02 & 0 & 2 & 0.167 & 0.537 & 14.175 & 0.483\\

\bottomrule\end{longtable}\endgroup

Исходные архивы траекторий содержат все независимые поступления и исходы кандидатных действий. Записи по обучающим наборам включают фактические пороги, их отношение к диапазону калибровки, активацию, время предупреждения, ложные и ненужные активации, интеграл WIP, длительности завершённых экземпляров и остаточную работу. Длительности завершённых экземпляров ограничены горизонтом оценки; проверки полной когорты, где доступны, являются отдельной диагностикой с продлённым горизонтом. Записи P4 сохраняют результаты обоих конкурирующих процессов.

\clearpage

\subsection{Диапазон наблюдений, активация и восстановление}

Таблица при нагрузочных возвратах показывает выбор retained-load. Числа порогов обозначают фактические значения WIP и число обучений. Диапазон обучения — разброс максимального наблюдавшегося WIP между наборами; $q/\max w$ считается внутри каждого набора. K/D обозначают известную/обнаруженную структуру. Диагностика не меняет правила активации.

\begin{center}\small\setlength{\tabcolsep}{3pt}\begin{tabular}{llp{.24\textwidth}rrr}\toprule

Process & N & Threshold:count & Train max & $q/\max w$ & Above\\\midrule

P1 & K & 6:9; 12:3 & 5--9 & 0.67--2.00 & 6/12\\

P1 & D & 6:9; 12:3 & 5--9 & 0.67--2.00 & 6/12\\

P2 & K & 8:9; 16:3 & 6--13 & 0.62--2.67 & 6/12\\

P2 & D & 8:9; 16:3 & 6--13 & 0.62--2.67 & 6/12\\

P3 & K & 10:11; 40:1 & 7--11 & 0.91--5.00 & 7/12\\

P3 & D & 10:11; 40:1 & 7--11 & 0.91--5.00 & 7/12\\

P4 & K & 10:11; 20:1 & 10--21 & 0.48--1.43 & 1/12\\

P4 & D & 10:11; 20:1 & 10--21 & 0.48--1.43 & 1/12\\

P5 & K & 10:12 & 11--18 & 0.56--0.91 & 0/12\\

P5 & D & 10:12 & 11--18 & 0.56--0.91 & 0/12\\

P6 & K & 10:11; 20:1 & 10--16 & 0.62--2.00 & 1/12\\

P6 & D & 10:11; 20:1 & 10--16 & 0.62--2.00 & 1/12\\

\bottomrule\end{tabular}\end{center}

\begin{center}\small\setlength{\tabcolsep}{4pt}\begin{tabular}{llrrrrr}\toprule

Process & N & Activate & Nonsevere & Lead/$T$ & Event & Prevented\\\midrule

P1 & K & 0.705 & 0.342 & 5.901 & 0.365 & 0.329\\

P1 & D & 0.705 & 0.342 & 5.901 & 0.365 & 0.329\\

P2 & K & 0.705 & 0.490 & 5.042 & 0.345 & 0.179\\

P2 & D & 0.705 & 0.490 & 5.042 & 0.345 & 0.179\\

P3 & K & 0.677 & 0.031 & 4.247 & 0.642 & 0.036\\

P3 & D & 0.677 & 0.031 & 4.247 & 0.642 & 0.036\\

P4 & K & 0.812 & 0.506 & 7.837 & 0.495 & 0.200\\

P4 & D & 0.812 & 0.506 & 7.837 & 0.495 & 0.200\\

P5 & K & 0.868 & 0.720 & 9.730 & 0.382 & 0.274\\

P5 & D & 0.868 & 0.720 & 9.730 & 0.382 & 0.274\\

P6 & K & 0.762 & 0.320 & 6.301 & 0.382 & 0.412\\

P6 & D & 0.762 & 0.320 & 6.301 & 0.382 & 0.412\\

\bottomrule\end{tabular}\end{center}

Nonsevere — активация при отсутствии тяжёлого события в парной траектории без действия; это не обязательно вредное вмешательство. Время предупреждения измеряется относительно начала такого контрфактического события. Предотвращение условно по событию без действия. Восстановление ниже означает WIP не выше $2S$ в течение $10T$ после импульса с цензурированием на исходном горизонте $180T$. Сначала усредняются траектории каждого независимого обучения, затем сами обучения; учитываются все импульсные траектории, а не только кризисные.

\begin{center}\small\begin{tabular}{lrrrr}\toprule

Process & No action & Predictive K & Retained K & Generator\\\midrule

P1 & 0.771 & 1.000 & 1.000 & 1.000\\

P2 & 0.745 & 0.979 & 0.979 & 1.000\\

P3 & 0.195 & 0.659 & 0.599 & 0.732\\

P4 & 0.536 & 0.982 & 0.867 & 0.987\\

P5 & 0.802 & 0.865 & 0.943 & 1.000\\

P6 & 0.482 & 0.948 & 0.987 & 1.000\\

\bottomrule\end{tabular}\end{center}

\subsection{Последствия общих ресурсов и положение узкого места}

Ресурсная диагностика повторяет P4/P5 при всех четырёх механизмах и всех 12 обучениях для двух импульсных условий с высокой базовой интенсивностью и повтора 0. Включены все типы действий на минимальном фиксированном пороге и сохранённые выбранные политики. Все 884 префиксов до $180T$ точно совпадают с исходными тестовыми отсчётами. Затем поступления прекращаются, а действия сохраняются до $600T$ для завершения когорты. Интегралы очередей относятся только к исходному горизонту. Поэтому длительность полной когорты имеет другой горизонт, чем основной показатель.

\begin{longtable}{llrrrr}\caption{P4: интегралы WIP каждого процесса (тысячи) по полному тесту и времена завершения полной когорты по продлённой диагностике.}\\\toprule

Mechanism & Action at 10 & WIP X & WIP Y & Cycle X & Cycle Y\\\midrule\endfirsthead

\toprule Mechanism & Action at 10 & WIP X & WIP Y & Cycle X & Cycle Y\\\midrule\endhead

fixed & None & 7.64 & 8.41 & 23.27 & 25.61\\

fixed & Capacity & 4.62 & 5.31 & 13.21 & 15.03\\

fixed & Rework X & 7.33 & 8.26 & 22.19 & 25.05\\

fixed & Rework Y & 7.34 & 7.82 & 22.09 & 23.40\\

service & None & 43.69 & 45.71 & 176.41 & 181.78\\

service & Capacity & 22.70 & 24.06 & 75.40 & 78.33\\

service & Rework X & 41.75 & 44.77 & 165.43 & 175.66\\

service & Rework Y & 42.15 & 42.57 & 166.09 & 164.01\\

routing & None & 31.65 & 30.88 & 119.63 & 116.00\\

routing & Capacity & 15.28 & 15.16 & 51.39 & 50.36\\

routing & Rework X & 16.83 & 22.30 & 58.65 & 77.74\\

routing & Rework Y & 21.92 & 16.61 & 76.21 & 56.21\\

combined & None & 90.99 & 89.23 & 538.82 & 525.69\\

combined & Capacity & 71.85 & 70.13 & 347.75 & 338.59\\

combined & Rework X & 69.00 & 83.06 & 333.96 & 441.36\\

combined & Rework Y & 85.34 & 69.00 & 456.45 & 335.60\\

\bottomrule\end{longtable}

\begin{longtable}{llrrr}\caption{P5: средние очереди B и D на исходном горизонте; число парных траекторий с перемещением наибольшей очереди от B к D.}\\\toprule

Mechanism & Action at 10 & Queue B & Queue D & B to D\\\midrule\endfirsthead

\toprule Mechanism & Action at 10 & Queue B & Queue D & B to D\\\midrule\endhead

fixed & None & 12.13 & 0.11 & 0/24\\

fixed & Capacity B & 1.76 & 4.55 & 24/24\\

fixed & Capacity D & 12.13 & 0.01 & 0/24\\

fixed & Rework & 12.13 & 0.08 & 0/24\\

service & None & 140.59 & 0.01 & 0/24\\

service & Capacity B & 38.34 & 0.11 & 0/24\\

service & Capacity D & 140.59 & 0.00 & 0/24\\

service & Rework & 140.59 & 0.00 & 0/24\\

routing & None & 12.17 & 36.45 & 0/24\\

routing & Capacity B & 1.67 & 46.10 & 0/24\\

routing & Capacity D & 12.17 & 4.38 & 0/24\\

routing & Rework & 12.17 & 1.73 & 0/24\\

combined & None & 137.95 & 1.28 & 0/24\\

combined & Capacity B & 37.29 & 116.13 & 24/24\\

combined & Capacity D & 137.95 & 0.05 & 0/24\\

combined & Rework & 137.95 & 0.04 & 0/24\\

\bottomrule\end{longtable}


\label{maintextend}
\bibliographystyle{splncs04}
\bibliography{references}
\end{document}
